\documentclass[11pt]{article}
\usepackage[T1]{fontenc}
\usepackage{lmodern,microtype,amsmath,amssymb,amsthm,mathtools,mathrsfs,tikz-cd}
\usepackage[margin=1in]{geometry}
\usepackage{xurl,hyperref}
\hypersetup{colorlinks=true,linkcolor=blue,citecolor=blue,urlcolor=blue,
pdftitle={Local connectedness in recursion categories and essential incompleteness},pdfauthor={Florian Lengyel}}
\newtheorem{theorem}{Theorem}[section]
\newtheorem{proposition}[theorem]{Proposition}
\newtheorem{lemma}[theorem]{Lemma}
\newtheorem{corollary}[theorem]{Corollary}
\newcommand{\N}{\mathbb N}
\newcommand{\PA}{\mathrm{PA}}
\newcommand{\Pfn}{\mathbf{Pfn}}
\newcommand{\ran}{\operatorname{ran}}
\title{Local connectedness in recursion categories\\
and essential incompleteness}
\author{Florian Lengyel}
\date{Version 7\\31 October 2026}
\begin{document}
\maketitle

\begin{abstract}
We prove a range-functor criterion for recursive inseparability in a
recursion category, with local connectedness required only in the target.
The criterion assumes that the sections associated with complementary
source domains lift to source arrows; preservation of products is
unnecessary. A counterexample
shows that section lifting cannot be omitted, even for a faithful functor
preserving products and reflecting zero maps.

Miller's sequence yields a recursion category of partial numerical
functions computed independently of its binary name. The index sets of
terms below and above $1/2$ are inseparable in this category, although
every individual name computes a separator. A recursive reduction from
the categorical diagonal pair proves this assertion. The cut pair
gives essential incompleteness, relative to computation independent of
the name, for a theory of real closed fields with constants specifying
the sequence's rational cuts.

We also construct faithful model interpretations of the syntactic
recursion categories associated with consistent extensions of Peano
arithmetic. The criterion gives their inseparable diagonal pairs; its
numerical instance yields recursive inseparability of the theorem and
refutation sets of every consistent computably enumerable extension of
Peano arithmetic. A cosheaf construction gives a recursion category
not restriction-equivalent to the classical one and a consistent
propositional theory whose consistent
extensions represented by restriction idempotents are all incomplete.
\end{abstract}

\section{Introduction}

Theorem~\ref{thm:local-lifting} establishes recursive inseparability
in a recursion category using a range functor to a split locally
connected target satisfying weak choice and inverse-image preservation
of binary joins. The target provides sections associated with
complementary source restriction idempotents. If these sections lift
to source arrows, a Turing morphism gives an inseparable source pair.
The functor need not preserve restriction products.

Joins of source arrows with disjoint domains suffice for lifting
(Proposition~\ref{prop:disjoint-joins}). Without lifting the conclusion
can fail even for a faithful range functor preserving products and
reflecting zero maps
(Appendix~\ref{app:lifting-counterexample}).

The principal application uses a sequence $a:\N\to[0,1]$ constructed
by Miller \cite{Miller2004}. We form a category $\mathcal E_a$ whose
arrows are partial numerical functions computed by one oracle program
with the same domain and values for every binary name of $a$.
Compactness of the set of names permits the construction of a Turing
morphism in this category. The index sets
\[
 L_a=\{n:a(n)<1/2\},\qquad R_a=\{n:a(n)>1/2\}
\]
are semidecidable in $\mathcal E_a$. We construct an ordinary recursive
map reducing the diagonal pair of Theorem~\ref{thm:local-lifting} to
$L_a,R_a$, and hence prove their inseparability in $\mathcal E_a$.
For each fixed binary name $p$, however, the first digit of each row
computes a separator. The inclusion into the recursion category of
$p$-partial recursive functions is a faithful range functor preserving
products. A faithful range functor preserving products can therefore
send an inseparable pair to a separable one.

This example gives a theory of real closed fields with constants $c_n$
and the strict rational inequalities satisfied by $a(n)$.
The theory is consistent, and its axioms are semidecidable independently
of the name of $a$. Every consistent extension with the same
semidecidability property is incomplete
(Theorem~\ref{thm:miller-incompleteness}). The sentences $c_n<1/2$
pull the theorem and refutation domains back to supersets of $L_a$
and $R_a$, so Corollary~\ref{cor:domain-incompleteness} applies.

We also apply the criterion to the syntactic recursion categories of
Montagna and Di Paola--Montagna \cite{Montagna,DPM}.
For a consistent computably enumerable (c.e.)
extension $T$ of Peano arithmetic ($\PA$), the category $S_T$ identifies
programs when $T$ proves one universally quantified sentence asserting
equality of their partial graphs. These categories are not locally
connected \cite[Corollary 4.4]{DPM}. Interpretation in all models of
$T$ with carrier $\N$ gives a faithful range functor into a locally
connected product of partial-map categories. Programs selecting between
predicates proved complementary in $T$ provide the section lifts.
The construction commutes with theory extension.

The numerical instance of the criterion gives a fixed recursively
inseparable pair. Finite computation proofs show that its two sets lie
in the inverse images of the theorem and refutation sets of $\PA$ under
numeral substitution. The inclusions persist under theory extension.
Consequently every consistent c.e.\ extension of $\PA$ has recursively
inseparable theorem and refutation sets and is incomplete.
We also construct a recursion category from cosheaves on a three-point
space. It is not equivalent as a restriction category to the classical
recursion category. Its restriction splitting contains a consistent
propositional theory whose formula and derivation objects are
characterized by initiality. Every consistent extension represented
by a restriction idempotent is incomplete.

After the preliminaries, Section~\ref{sec:lifting} proves the general
criterion and its incompleteness corollary.
Section~\ref{sec:miller} gives the Miller application.
Sections~\ref{sec:syntactic}--\ref{sec:family} construct the syntactic
model functors and their section lifts. Sections~\ref{sec:arithmetic}
and~\ref{sec:cosheaves} give the arithmetic and cosheaf applications.
The appendices contain the cosheaf construction, its syntax and proofs,
and the counterexample to inseparability without section lifting.

\subsection{Related work}

Di Paola and Heller proposed inseparability of theorems and refutable
sentences as part of a categorical formulation of incompleteness
\cite[pp.~594--595]{DPH}. Local connectedness enters Heller's existence
theorem for recursion categories \cite[Sections 8--9]{Heller}.
Di Paola and Montagna's syntactic examples led to the question of
extending that theorem beyond locally connected sources \cite{DPM}.
Lengyel proved existence theorems using faithful iteration functors
\cite{Lengyel2004,Corrections}. The present criterion concerns
inseparability in a source already equipped with a Turing morphism.

Di Paola and Heller introduced dominical categories as a setting for
recursion theory \cite{DPH}. Rosolini's $P$-categories generalized this
setting \cite{RosoliniThesis,Rosolini}. Heller adopted them and defined
prodominical categories \cite[Section 5]{Heller}.
Cockett and Lack introduced restriction categories \cite{CockettLack},
and Cockett and Hofstra developed the theory of Turing categories
\cite{CockettHofstra}. We use restriction categories, with binary
restriction products in the source recursion categories. These products
give the equivalent $P$-category structure
\cite[Proposition 4.5]{CockettLackIII}.

Di Paola and Heller proved recursive inseparability in dominical
recursion categories satisfying the axiom of choice
\cite[Theorem 8.15]{DPH}. Bernardi obtained categorical inseparability
results from precomplete equivalence relations, including PA-provable
program equality among his examples \cite[pp.~125--128]{Bernardi}.
Montagna also obtained effectively inseparable pairs in $S$
\cite[Remark 3.1]{Montagna}.


Miller constructed a sequence that enumerates every real in $[0,1]$
computable uniformly from its representations \cite[Theorem~6.3]{Miller2004}.
We use the stronger index property stated immediately after that
theorem and recalled in Section~\ref{sec:miller}.
Bauer and Hanson use this construction in a parameterized realizability
topos with countable Dedekind reals \cite{BauerHanson}.
Section~\ref{sec:miller} uses their binary-name formulation to construct
a recursion category, identify an inseparable pair of strict-cut
predicates, and apply it to a theory of real closed fields.

\section{Categorical preliminaries}\label{sec:preliminaries}

\subsection{Partial maps and restriction idempotents}

Composition is written $gf=g\circ f$, so $f$ is applied first.
A \emph{restriction category} is a category equipped with an assignment
that sends every arrow
$f:X\rightharpoonup Y$ to an arrow
$\bar f:X\rightharpoonup X$ satisfying
\[
 f\bar f=f,\qquad
 \bar f\bar g=\bar g\bar f,\qquad
 \overline{g\bar f}=\bar g\bar f,\qquad
 \bar h f=f\overline{hf}.
\]
In the middle two equations $g:X\rightharpoonup Z$; in the last,
$h:Y\rightharpoonup Z$ \cite{CockettLack}.
The arrows $\bar f$ are idempotent and are called
\emph{restriction idempotents}.
The arrow $f$ is \emph{total} when $\bar f=1_X$.
The total arrows form a subcategory, denoted by
$\operatorname{Tot}(\mathcal C)$ for a restriction category $\mathcal C$.
A \emph{restriction functor} $F:\mathcal C\to\mathcal D$ between
restriction categories satisfies $F(\bar f)=\overline{F(f)}$ for every
arrow $f$. A \emph{restriction equivalence} consists of restriction
functors in both directions whose composites are naturally isomorphic
to the respective identity functors.

Restriction idempotents on $X$ are ordered by $a\le b$ when $ba=a$.
Their meet is composition, their greatest element is $1_X$, and
$a\vee b$ denotes their least upper bound when it exists.
For parallel partial arrows, the restriction order is
$f\le g$ when $f=g\bar f$.
Parallel arrows $f,g:X\rightharpoonup Y$ are \emph{compatible} when
$f\bar g=g\bar f$. Their join, when it exists, is their least upper
bound in the restriction order.
A \emph{restriction zero} $0_{X,Y}:X\rightharpoonup Y$ is
absorbing under composition and satisfies
$\overline{0_{X,Y}}=0_{X,X}$. We abbreviate $0_{X,X}$ to $0_X$.
Restriction idempotents $a,b$ on $X$ are disjoint when $ab=0_X$.
They are complementary when also $a\vee b=1_X$.
A disjoint pair $a_0,a_1$ is \emph{inseparable} if no complementary
pair $c,c'$ satisfies $a_0\le c$ and $a_1\le c'$.

For $f:X\rightharpoonup Y$ and a restriction idempotent $a$ on $Y$,
the \emph{inverse image} of $a$ is the restriction idempotent
$f^*a=\overline{af}$ on $X$.
A restriction idempotent $r$ on $Y$ \emph{receives} $f$ when $rf=f$
\cite[Section 6]{Heller}.
The \emph{range} $\ran(f)$ is the least restriction idempotent receiving $f$;
we require the agreement law
\[
 gf=hf\quad\Longrightarrow\quad
 g\ran(f)=h\ran(f)
 \qquad(g,h:Y\rightharpoonup Z).
\]
For an ordinary partial function $f$, the restriction idempotent
$\bar f$ and the range $\ran(f)$ are the partial identities on its
domain of definition and image, respectively.

A restriction idempotent $e$ on $X$ \emph{splits} if
$e=mr$ for a total $m:A\to X$ and $r:X\rightharpoonup A$
with $rm=1_A$. The category is \emph{split} when every restriction
idempotent splits.
Its \emph{restriction splitting}, denoted by
$\operatorname{Split}(\mathcal C)$, has objects $(X,e)$, where $e$ is a
restriction idempotent on $X$, and arrows
$f:(X,e)\rightharpoonup(Y,d)$ given by original arrows
$f:X\rightharpoonup Y$ satisfying $f=dfe$.
The identity of $(X,e)$ is $e$, and $f$ is total there when
$\bar f=e$. Thus $(X,e)$ is the object through which $e$ splits.

\subsection{Products, coproducts, and local connectedness}

A \emph{binary restriction product} $X\times Y$ has total
projections $\pi_1:X\times Y\to X$ and $\pi_2:X\times Y\to Y$.
For arrows $f:Z\rightharpoonup X$ and $g:Z\rightharpoonup Y$,
there is a unique $\langle f,g\rangle:Z\rightharpoonup X\times Y$
with
\[
 \pi_1\langle f,g\rangle=f\bar g,\qquad
 \pi_2\langle f,g\rangle=g\bar f,\qquad
 \overline{\langle f,g\rangle}=\bar f\bar g.
\]
Products of arrows are obtained by this pairing and the projections.
The total categorical diagonal is
$\Delta_X=\langle1_X,1_X\rangle:X\to X\times X$.
A restriction category has a $P$-category structure if and only if
it has binary restriction products
\cite[Proposition 4.5]{CockettLackIII}.

Binary \emph{restriction coproducts} are categorical coproducts with
total injections $i_0:X\to X+Y$ and $i_1:Y\to X+Y$.
For arrows $f:X\rightharpoonup Z$ and $g:Y\rightharpoonup Z$,
write $[f,g]:X+Y\rightharpoonup Z$ for their copairing.
The \emph{codiagonal} is
$\nabla_X=[1_X,1_X]:X+X\to X$.
In a category with ranges and restriction coproducts,
$a\vee b=\ran[a,b]$ for restriction idempotents on the same object.

A \emph{range functor} $F:\mathcal C\to\mathcal D$ preserves
restrictions, ranges, restriction zeros, and binary
restriction coproducts. In particular, the canonical map
\[
 \chi_{X,Y}=[F i_0,F i_1]:FX+FY\longrightarrow F(X+Y)
\]
is an isomorphism, and
$F[f,g]\chi_{X,Y}=[Ff,Fg]$ for $f:X\rightharpoonup Z$ and
$g:Y\rightharpoonup Z$. A range functor preserves binary
joins of restriction idempotents by the preceding formula, and hence
complementary pairs.
It need not preserve restriction products. Let $\Pfn$ denote the
category of sets and partial functions. For example, the functor
$F:\Pfn\to\Pfn$ given by $F(X)=X\times\{0,1\}$ and
$F(f)=f\times1_{\{0,1\}}$ preserves restrictions, images, zeros,
and coproducts. Its canonical map from $F(X\times Y)$ to
$FX\times FY$ sends $((x,y),a)$ to $((x,a),(y,a))$, and is
not surjective when $X,Y$ are nonempty.

A restriction idempotent $a$ on $X$ is \emph{connected} if, for
restriction idempotents $b,c$ on $X$, the equations $a=b\vee c$ and
$bc=0$ imply $b=0$ or $c=0$. The category is \emph{locally connected}
when its connected restriction idempotents are jointly epimorphic: for
$f,g:X\rightharpoonup Y$, the equations $fa=ga$ for every connected
restriction idempotent $a$ on $X$ imply $f=g$.
A \emph{partial monomorphism} $p:X\rightharpoonup Y$ satisfies
$pu=pv\Rightarrow\bar p u=\bar p v$ for all
$u,v:W\rightharpoonup X$.
A \emph{section} of $p$ is an arrow $s:Y\rightharpoonup X$
with $ps=\bar s$ and $psp=p$.
\emph{Weak choice} means that every partial monomorphism has such
a section. For partial functions, this is the existence of the
partial inverse of an injection on its domain.

\subsection{Recursion categories}

A restriction category with binary restriction products and zeros is
\emph{prodominical} if $f\times0=0$ for every arrow $f$ and every
zero arrow of the required type. It is \emph{isotypical} if any
two of its objects are isomorphic.

A \emph{Turing morphism} on $U$ is an arrow
$\tau:U\times U\rightharpoonup U$ such that, for every object $X$
and every $f:X\times U\rightharpoonup U$, there is a total
parameter map $r:X\to U$ satisfying
\[
 f=\tau(r\times1_U).
\]
A \emph{recursion category} here is a prodominical isotypical
category with binary restriction products and a Turing morphism,
as in \cite[Section 9]{Heller}.

\section{Switching and the local-connectedness theorem}\label{sec:lifting}

\subsection{The local-connectedness construction}

For a coproduct $Y+Z$, let
$\varepsilon_0=1_Y+0_Z$ and $\varepsilon_1=0_Y+1_Z$ be its summand
restriction idempotents. Their meet is zero and their join is the identity. These facts
follow by precomposition with the two injections and the agreement
property of ranges.

\begin{lemma}\label{lem:connected-alternative}
Let $\mathcal A$ be a restriction category with ranges, restriction zeros,
and binary restriction coproducts. Suppose inverse images preserve binary
joins of restriction idempotents. If $a$ is a connected restriction
idempotent on $X$ and
$f:X\to Y+Z$ is total, then
$\varepsilon_0fa=fa$ or $\varepsilon_1fa=fa$.
\end{lemma}
\begin{proof}
Totality and the inverse-image laws give the disjoint decomposition
\[
 a=af^*\varepsilon_0\vee af^*\varepsilon_1.
\]
Connectedness makes one term zero and the other equal to $a$. In the
latter case $a\le f^*\varepsilon_i$, and the restriction identity
$\varepsilon_i f=f\overline{\varepsilon_i f}$ gives
$\varepsilon_i fa=fa$.
\end{proof}

\begin{proposition}\label{prop:local-switch}
Let $\mathcal A$ be a locally connected split restriction category with
ranges, restriction zeros, and binary restriction coproducts. Suppose
weak choice holds and inverse images preserve binary joins of
restriction idempotents.
For disjoint restriction idempotents $a,b$ on $X$, the arrow
$p=\nabla_X(a+b):X+X\rightharpoonup X$ is a partial monomorphism
with a unique section $s:X\rightharpoonup X+X$, and
\[
 ps=\bar s=a\vee b.
\]
In particular, for complementary $c,c'$ on $X$, the arrow
$\nabla_X(c'+c)$ has a total section $s:X\to X+X$ satisfying
$sc=i_1c$ and $sc'=i_0c'$.
\end{proposition}
\begin{proof}
First let $\theta,\theta':Y\to X+X$ be total and suppose
$p\theta=p\theta'$. The identities
$ap=\nabla_X(a+0)$ and $bp=\nabla_X(0+b)$ give
\[
 \nabla_X(a+0)\theta=\nabla_X(a+0)\theta',
 \qquad
 \nabla_X(0+b)\theta=\nabla_X(0+b)\theta'.
\]
For a connected restriction idempotent $\gamma$ on $Y$,
Lemma~\ref{lem:connected-alternative} places each of
$\theta\gamma,\theta'\gamma$ in one summand.
If both lie in the first summand, the first equality, followed by $i_0$,
gives $(a+b)\theta\gamma=(a+b)\theta'\gamma$; the second-summand
case uses the second equality and $i_1$.
If they lie in opposite summands, the two equalities give
$(a+b)\theta\gamma=0=(a+b)\theta'\gamma$.
Local connectedness therefore gives
$(a+b)\theta=(a+b)\theta'$. Since $\bar p=a+b$, this proves the
partial-monomorphism condition for total arrows.

For arbitrary $u,v:W\rightharpoonup X+X$ with $pu=pv$, split their
common restriction idempotent as
$\overline{pu}=\overline{pv}=me$, where $m:Z\to W$ is total,
$e:W\rightharpoonup Z$, and $em=1_Z$. The arrows $pum,pvm$, hence also $um,vm$, are total.
The preceding case gives $\bar p um=\bar p vm$. Composing with $e$
and using $\bar p u=\bar p u\,\overline{pu}$ and its analogue for $v$
gives $\bar p u=\bar p v$. Thus $p$ is a partial monomorphism.

Weak choice supplies a section $s$. The equations $ps=\bar s$ and
$psp=p$ imply $\bar s=\ran(p)$: the idempotent $\bar s$ receives $p$,
and every restriction idempotent receiving $p$ also receives
$ps=\bar s$. Moreover
$\bar p s=s\overline{ps}=s$. Two sections have the same composite
with $p$, so partial monicity makes them equal.
The defining property of the range gives $\ran(p)=a\vee b$.

For $a=c'$ and $b=c$, the section is total because $c'\vee c=1_X$.
The equality $p(sc)=c=p(i_1c)$ and partial monicity give $sc=i_1c$;
the equation on $c'$ is analogous.
\end{proof}

\begin{theorem}[Range-functor lifting]\label{thm:local-lifting}
Let $\mathscr R$ be a recursion category with ranges and
binary restriction coproducts. Let $\mathcal A$ be a locally connected
split restriction category with ranges, restriction zeros, and binary
restriction coproducts, satisfying weak choice, with inverse images
preserving binary joins of restriction idempotents. Let
$F:\mathscr R\to\mathcal A$ be a range functor.
Choose a Turing morphism $\tau:U\times U\rightharpoonup U$ and
an isomorphism $b:U\to U+U$, with $F(1_U)\ne0$.
For each complementary pair $c,c'$ on $U$, let
$s_c:FU\to FU+FU$ be the total section of $\nabla_{FU}(Fc'+Fc)$ supplied by
Proposition~\ref{prop:local-switch}. Suppose there is an arrow
$\ell_c:U\rightharpoonup U+U$ with
$F\ell_c=\chi_{U,U}s_c$, where
$\chi_{U,U}:FU+FU\to F(U+U)$ is the canonical coproduct isomorphism.
Then
\[
 a_0=(b\tau\Delta_U)^*(1_U+0_U),\qquad
 a_1=(b\tau\Delta_U)^*(0_U+1_U)
\]
are inseparable in $\mathscr R$.
\end{theorem}
\begin{proof}
The restriction idempotents $a_0,a_1$ are disjoint. Suppose $c,c'$ separate them.
Put $u=b^{-1}\ell_c$, and choose a total parameter $g:U\to U$ with
$\tau(g\times1_U)=u\pi_2$. Precomposition with $\langle1_U,g\rangle$
gives $\tau\Delta_Ug=ug$.

Applying $F$ to the source equation
$b\tau\Delta_Ug=\ell_cg$ gives
\[
 \chi_{U,U}^{-1}F(b\tau\Delta_U)Fg=s_cFg.
\]
In the target write $E=Fg:FU\to FU$,
$H=\chi_{U,U}^{-1}F(b\tau\Delta_U):FU\rightharpoonup FU+FU$,
$\Gamma=Fc$, and $\Gamma'=Fc'$. Thus $HE=s_cE$.
Since $s_c\Gamma=i_1\Gamma$ and $s_c\Gamma'=i_0\Gamma'$,
the inverse images under $s_c$ of the two summand restriction idempotents are
$\Gamma'$ and $\Gamma$,
respectively. Preservation of restrictions and composition now gives
\[
 E^*F(a_0)=E^*\Gamma',\qquad E^*F(a_1)=E^*\Gamma.
\]
The separator inequalities imply that $E^*\Gamma$ and $E^*\Gamma'$
are equal. They are complementary because $E$ is total and target
inverse images preserve joins. Hence $1_{FU}=0$, contradicting the
hypothesis.
\end{proof}

\subsection{A sufficient condition for lifting}

\begin{proposition}\label{prop:disjoint-joins}
Let $\mathcal C$ be a restriction category with restriction zeros and
binary restriction coproducts. Suppose every pair of parallel arrows
whose restriction idempotents are disjoint has a join in the restriction
order.
For complementary restriction idempotents $c,c'$ on $X$, the arrow
\[
 \ell_c=(i_1c)\vee(i_0c'):X\rightharpoonup X+X
\]
is total and satisfies $\ell_cc=i_1c$ and $\ell_cc'=i_0c'$.
If $F:\mathcal C\to\mathcal A$ is a range functor and $\mathcal A$
satisfies the hypotheses of Proposition~\ref{prop:local-switch},
then $\chi_{X,X}^{-1}F\ell_c$ is the section of
$\nabla_{FX}(Fc'+Fc)$.
\end{proposition}
\begin{proof}
The restrictions of $i_1c$ and $i_0c'$ are $c$ and $c'$ because
the injections are total. The upper-bound equations for the join
give $\ell_cc=i_1c$ and $\ell_cc'=i_0c'$.
Both $c$ and $c'$ lie below $\bar\ell_c$, so
$1_X=c\vee c'\le\bar\ell_c$ and $\ell_c$ is total.

Put $\Gamma=Fc$ and $\Gamma'=Fc'$. Preservation of composition
and coproduct injections gives
$(\chi_{X,X}^{-1}F\ell_c)\Gamma=i_1\Gamma$ and
$(\chi_{X,X}^{-1}F\ell_c)\Gamma'=i_0\Gamma'$.
The section $s$ of Proposition~\ref{prop:local-switch} satisfies the
same two equations. Thus $\chi_{X,X}^{-1}F\ell_c$ and $s$ agree after
$[\Gamma,\Gamma']$. Its range is
$\Gamma\vee\Gamma'=1_{FX}$, so range agreement gives
$\chi_{X,X}^{-1}F\ell_c=s$.
\end{proof}

\subsection{From inseparable restriction idempotents to incompleteness}

Inseparability of theorem and refutation domains follows when their
inverse images under a total sentence map contain an inseparable pair.

\begin{corollary}\label{cor:domain-incompleteness}
Let $\mathcal C$ be a restriction category with restriction zeros and
binary joins of restriction idempotents preserved by inverse image along total arrows.
Let $a,b$ be inseparable restriction idempotents on $U$. Suppose that
$\theta:U\to S$ and $\nu:S\to S$ are total, and that a restriction
idempotent $t_0$ on $S$ satisfies
\[
 a\le\theta^*t_0,\qquad b\le\theta^*\nu^*t_0.
\]
For every restriction idempotent $t\ge t_0$ on $S$ with $t\,\nu^*t=0_S$,
the pair $t,\nu^*t$ is inseparable. In particular,
$t\vee\nu^*t\ne1_S$.
\end{corollary}
\begin{proof}
Suppose complementary restriction idempotents $e,e'$ on $S$ satisfy
$t\le e$ and $\nu^*t\le e'$. Put
$c=\theta^*e$ and $c'=\theta^*e'$. Inverse image preserves
meets and zero, so $cc'=\theta^*(ee')=0_U$.
Preservation of joins and totality of $\theta$ give
\[
 c\vee c'=\theta^*(e\vee e')=\theta^*1_S=1_U.
\]
Monotonicity gives
$a\le\theta^*t_0\le\theta^*t\le c$ and
$b\le\theta^*\nu^*t_0\le\theta^*\nu^*t\le c'$.
Thus $c,c'$ would separate $a,b$, a contradiction.
If $t\vee\nu^*t=1_S$, the disjoint pair $t,\nu^*t$ would
itself be such a complementary separator.
\end{proof}

For a sentence object $S$, a negation map $\nu$, and a restriction
idempotent $t$ representing the theorems, the equations
$t\,\nu^*t=0_S$ and $t\vee\nu^*t=1_S$
express consistency and completeness, respectively.
If extensions of the base theory are represented by restriction
idempotents $t\ge t_0$ and the two inequalities hold for the base
theory, the corollary rules out a complete consistent extension
in this class.

An inseparable pair on an object of $\mathscr R$ remains inseparable
in $\operatorname{Split}(\mathscr R)$: the fully faithful inclusion
$U\mapsto(U,1_U)$ changes neither its endomorphisms nor its restriction idempotents.
Thus sentence and proof objects may be constructed in the restriction
splitting when applying Corollary~\ref{cor:domain-incompleteness}.

\section{Miller sequences and essential incompleteness}\label{sec:miller}

Fix a recursive bijection $\langle-,-\rangle:\N^2\to\N$ and an
acceptable enumeration $(\Phi_e^p)_{e\in\N}$ of partial numerical
oracle computations, where $p\in2^{\N}$. Tuple inputs and outputs use
fixed recursive codings. Let $(q_j)_{j\in\N}$ be a recursive enumeration
of the rationals. A program $d$ \emph{computes a real $r$ from $p$}
if $\Phi_d^p(k)$ is defined for every $k$ and
\[
 |q_{\Phi_d^p(k)}-r|<2^{-k}\qquad(k\in\N).
\]
Put $Q=[0,1]^{\N}$ and define
\[
 \rho:2^{\N}\longrightarrow Q,\qquad
 \rho(p)(n)=\sum_{j\geq0}p(\langle n,j\rangle)2^{-j-1}.
\]
A \emph{binary name} of $z\in Q$ is an element of
$P_z=\rho^{-1}(\{z\})$. The map $\rho$ is continuous and onto,
so every $P_z$ is nonempty and compact.

Miller's construction gives a sequence $a\in Q$ with the following
property: if a program $d$ computes the same real $r\in[0,1]$ from
every $p\in P_a$, then
\begin{equation}\label{eq:miller-property}
 a(d)=r.
\end{equation}
The index equality is the strengthened conclusion stated after
\cite[Theorem~6.3]{Miller2004}; we use the binary-name formulation in
\cite[Section~2.1 and Theorem~2.2]{BauerHanson}.
The rational approximations produced by $d$ may depend on the name,
while the real they represent is fixed.

\subsection{Computation independent of the binary name}
\label{sec:miller-category}

\begin{lemma}\label{lem:miller-unanimity}
There is one oracle program which, from $p\in P_z$ and numerical inputs
$e,n$, computes the partial function
\[
 \tau_z(e,n)=m
 \quad\Longleftrightarrow\quad
 (\forall q\in P_z)\;\Phi_e^q(n)\mathord\downarrow=m.
\]
Its value and its domain of definition depend only on $z$, not on the
chosen name $p$.
\end{lemma}
\begin{proof}
For a finite binary word $\sigma$, write $[\sigma]$ for its cylinder in
Cantor space. The set $\rho([\sigma])$ is a computably specified finite
union of closed dyadic boxes in $Q$, with restrictions on only finitely
many coordinates. Indeed, the finitely many prescribed bits can be
extended in finitely many ways to initial segments of each affected row.
From a name $p$ of $z$, strict rational inequalities for the coordinates
of $z$ are semidecidable. We can therefore enumerate all cylinders
$[\sigma]$ for which $z\notin\rho([\sigma])$. Their union is exactly
$2^{\N}\setminus P_z$: if $\rho(q)\ne z$, a sufficiently long initial
segment of $q$ already forces a different value in one coordinate.

For fixed $e,n,m$, let
\[
 H_{e,n,m}=\{q\in2^{\N}:\Phi_e^q(n)\mathord\downarrow=m\}.
\]
Finite halting computations enumerate this open set as a union of
cylinders, uniformly in $e,n,m$. For each $m$, search for a finite cover
of Cantor space by cylinders from $H_{e,n,m}$ and
$2^{\N}\setminus P_z$, dovetailing these searches over $m$. Whether finitely
many cylinders cover Cantor space is decidable. Such a cover exists
exactly when $P_z\subseteq H_{e,n,m}$, by compactness. Since $P_z$ is
nonempty, at most one $m$ satisfies this condition. Thus the search has
the asserted output and definedness, for every name of $z$.
\end{proof}

For the fixed sequence $a$, let $\mathcal E_a$ have the
objects generated from $\N$ by binary Cartesian products and tagged
disjoint unions, with fixed recursive codings. Every such object is
recursively isomorphic to $\N$. An arrow $f:X\rightharpoonup Y$ belongs
to $\mathcal E_a$ when one oracle program computes exactly $f$, including
its domain of definition, for every $p\in P_a$. Equality of arrows is
equality of partial functions.

\begin{proposition}\label{prop:miller-category}
The category $\mathcal E_a$ is a locally connected recursion category
with ranges, binary restriction coproducts, and finite compatible joins.
It satisfies weak choice, and $\tau_a:\N\times\N\rightharpoonup\N$
is a Turing morphism.
\end{proposition}
\begin{proof}
Composition of the representing programs gives the same composite
partial function for every name. The recursive projections, pairings,
tagged injections, and case distinctions give the stated finite
structure. Restriction is computed by running a program for $f(x)$ and,
if it halts, returning $x$. To compute the range partial identity at
$y$, dovetail the computations $f(x)$ over all coded inputs $x$ and
return $y$ when one returns $y$. Although the discovered witness may
depend on the name, the returned value and the existence of a witness
do not. Thus $\ran(f)$ belongs to $\mathcal E_a$; its leastness and
agreement law follow from those of set-theoretic images.
Dovetailing compatible computations gives their union;
compatibility ensures that the output is independent of which
computation finishes first. These operations also give zero maps and
the prodominical law $f\times0=0$.

For every input $x$, the singleton partial identity $1_{\{x\}}$ is
recursive. It is connected, and these partial identities detect
inequality of partial functions, so the category is locally connected.
A partial monomorphism here is injective on its domain: a failure of
injectivity would be detected by two recursive constant maps. For an
injective arrow $f$, the same preimage search returns the unique $x$
with $f(x)=y$. Its output is consequently independent of the name and
gives the partial inverse. This proves weak choice.

Lemma~\ref{lem:miller-unanimity} puts $\tau_a$ in $\mathcal E_a$.
Let $f:X\times\N\rightharpoonup\N$ be represented by a program with
index $d$. The parameter theorem for oracle programs gives an ordinary
total recursive map $c:X\to\N$ with
\[
 \Phi_{c(x)}^p(n)\simeq\Phi_d^p(x,n).
\]
Here $\simeq$ includes simultaneous undefinedness.
If $f(x,n)=m$, every name gives this same output; if $f(x,n)$ is
undefined, every name gives an undefined computation. The definition
of $\tau_a$ therefore gives
$f=\tau_a(c\times1_{\N})$. The map $c$ is a total source arrow.
Finally, the chosen recursive codings make the category isotypical.
\end{proof}

For $p\in P_a$, let $\mathrm{Rec}^{p}$ be the recursion category on
the same numerical objects, with partial $p$-recursive functions as
arrows. The identity interpretation of partial functions defines a
faithful range functor
\[
 J_p:\mathcal E_a\longrightarrow\mathrm{Rec}^{p}.
\]
It also preserves restriction products and finite compatible joins.
Write $\mathcal K_p=\operatorname{Split}(\mathrm{Rec}^{p})$ and let
$I_p:\mathrm{Rec}^{p}\to\mathcal K_p$ be the full inclusion.
A concrete presentation of $\mathcal K_p$ has $p$-c.e. subsets of the
numerical objects as objects and partial $p$-recursive functions between
these carriers as arrows. A restriction idempotent splits through its $p$-c.e. domain.
Singleton restrictions prove local connectedness, and preimage search
for an injective map proves weak choice. Ranges are images, and inverse
images preserve domain unions. Thus $\mathcal K_p$ satisfies the target
hypotheses of Theorem~\ref{thm:local-lifting}.

If complementary source domains $c,c'$ partition a numerical object
$X$, run their two membership computations in parallel and return the
opposite tag when one succeeds. This gives the total source map
\[
 \ell_c=(i_1c)\vee(i_0c'):X\longrightarrow X+X.
\]
Its image under $I_pJ_p$ is the required target switch, and
$(I_pJ_p)(1_{\N})\ne0$. Thus each composite
$I_pJ_p$ satisfies the hypotheses of Theorem~\ref{thm:local-lifting}.
Using the recursive bijection $b:\N\to\N+\N$ which assigns the tag
of $m$ according to its parity, that theorem gives the inseparable
source domains whose supports are
\begin{equation}\label{eq:miller-diagonal-pair}
 A_i=\{e:\tau_a(e,e)\mathord\downarrow=m
                   \text{ for some }m\equiv i\pmod2\},
 \qquad i=0,1.
\end{equation}

\subsection{Strict cuts of the Miller sequence}
\label{sec:miller-cuts}

Define
\[
 L_a=\{n:a(n)<\tfrac12\},\qquad
 R_a=\{n:a(n)>\tfrac12\}.
\]
Their partial identities belong to $\mathcal E_a$, since strict
inequalities can be recognized from approximations obtained from any
binary name.

\begin{lemma}\label{lem:miller-cut-reduction}
There is an ordinary total recursive map $h:\N\to\N$ such that
\[
 h^{-1}(L_a)=A_0,\qquad h^{-1}(R_a)=A_1.
\]
\end{lemma}
\begin{proof}
For $e,m,t\in\N$, let $H_{e,m,t}$ be the finite union of halting
cylinders enumerated by stage $t$ for the computation
$\Phi_e^q(e)\mathord\downarrow=m$. Choose these finite unions increasing
in $t$, and put
\[
 K_{e,m,t}=\rho(2^{\N}\setminus H_{e,m,t})\subseteq Q.
\]
Each $K_{e,m,t}$ is a computably specified finite union of closed dyadic
boxes. Compactness of $P_z$ gives
\begin{align*}
 U_e^i
 &:=\{z:\tau_z(e,e)\mathord\downarrow=m
                      \text{ for some }m\equiv i\pmod2\}\\
 &=\bigcup_{\substack{m\equiv i\pmod2\\t\in\N}}
       (Q\setminus K_{e,m,t}).
\end{align*}
In particular, $U_e^0$ and $U_e^1$ are disjoint effectively open sets.

Equip $Q$ with the metric
\[
 d(z,w)=\sum_{n\geq0}2^{-n-1}|z(n)-w(n)|.
\]
For a nonempty box restricting finitely many coordinates to closed
rational intervals $I_n$, the distance from $z$ to the box is the
finite sum
$\sum_n2^{-n-1}\operatorname{dist}(z(n),I_n)$ over those coordinates.
Distance to a finite union is the minimum of these distances. Define
$\delta_{e,m,t}(z)=d(z,K_{e,m,t})$ when $K_{e,m,t}$ is nonempty, and
$\delta_{e,m,t}(z)=1$ when it is empty. These are uniformly computable
continuous functions with values in $[0,1]$, and
\[
 \delta_{e,m,t}(z)>0\quad\Longleftrightarrow\quad
 z\notin K_{e,m,t}.
\]
For each $i$, take a fixed recursive enumeration
$(m_j,t_j)_{j\geq0}$ of all pairs with $m_j\equiv i\pmod2$, and set
\[
 g_e^i(z)=\sum_{j\geq0}2^{-j-1}\delta_{e,m_j,t_j}(z).
\]
The tail bound is uniform, so $g_e^i$ is computable and continuous,
with $0\leq g_e^i\leq1$ and $g_e^i(z)>0$ exactly on $U_e^i$.
Consequently the computable continuous function
\[
 r_e(z)=\frac12+\frac{g_e^1(z)-g_e^0(z)}4
\]
takes values in $[1/4,3/4]$, is below $1/2$ exactly on $U_e^0$, and
is above $1/2$ exactly on $U_e^1$.

Uniform compilation supplies a total recursive $h$ such that program
$h(e)$, with oracle $p$, computes a rapidly converging rational name
of $r_e(\rho(p))$. For all $p\in P_a$ this is the same real $r_e(a)$.
The defining property of the Miller sequence therefore gives
\[
 a(h(e))=r_e(a).
\]
Together with $a\in U_e^i\Longleftrightarrow e\in A_i$, this proves
the asserted inverse-image equalities.
\end{proof}

\begin{corollary}\label{cor:miller-cuts-inseparable}
The domains $1_{L_a},1_{R_a}:\N\rightharpoonup\N$ are inseparable in
$\mathcal E_a$. Nevertheless, every $p\in P_a$ computes a separator
of the underlying sets $L_a,R_a$.
\end{corollary}
\begin{proof}
A complementary source separator for $L_a,R_a$ would pull back along
the total source map $h$ to a separator for $A_0,A_1$, contrary to
Theorem~\ref{thm:local-lifting} as applied above.
For a fixed name $p$, the set
\[
 C_p=\{n:p(\langle n,0\rangle)=0\}
\]
is $p$-decidable and satisfies $L_a\subseteq C_p$ and
$C_p\cap R_a=\varnothing$. At $a(n)=1/2$, either first binary digit
is allowed. Thus the single program reading that digit produces a
separator from every name, but its output set cannot be independent
of the name.
\end{proof}

In particular, $J_p\tau_a$ is not a Turing morphism in the target.
Indeed, $h^{-1}(C_p)$ is a $p$-decidable separator of its diagonal pair;
the diagonal argument for a Turing morphism would forbid such a
separator. The category $\mathrm{Rec}^{p}$ has its own Turing morphism.


\subsection{A theory of real closed fields}

Call a subset $D$ of a numerical object \emph{$\mathcal E_a$-semidecidable}
when its partial identity $1_D$ is an arrow of $\mathcal E_a$.
Consider the language of ordered fields with additional constants
$(c_n)_{n\in\N}$. Let $T_a$ consist of a recursive axiomatization of
real closed ordered fields together with the axioms
\begin{equation}\label{eq:miller-theory}
 q<c_n\quad\text{if }q<a(n),\qquad
 c_n<q\quad\text{if }a(n)<q,
 \qquad(q\in\mathbb Q, n\in\N).
\end{equation}
Rational bounds abbreviate inequalities with integer coefficients.
The real field, with $c_n$ interpreted as $a(n)$, is a model of $T_a$.
Its axiom set is $\mathcal E_a$-semidecidable: the field axioms are
recursive, and each strict rational inequality can be recognized from
every binary name of $a$ with the same membership outcome.

\begin{theorem}\label{thm:miller-incompleteness}
Let $a$ be a Miller sequence with property~\eqref{eq:miller-property},
and let $T_a$ be the theory~\eqref{eq:miller-theory}.
Every consistent extension of $T_a$ with an
$\mathcal E_a$-semidecidable axiom set has inseparable theorem and
refutation domains in $\operatorname{Split}(\mathcal E_a)$.
In particular, every such extension is incomplete.
\end{theorem}
\begin{proof}
The set $S$ of sentence codes is recursive. It therefore defines an
object of $\operatorname{Split}(\mathcal E_a)$. Negation and insertion
of a constant into an inequality give ordinary recursive total arrows
\[
 \nu:S\to S,\qquad
 \nu(\ulcorner\varphi\urcorner)=\ulcorner\neg\varphi\urcorner,
 \qquad
 \theta:\N\to S,\qquad
 \theta(n)=\ulcorner c_n<1/2\urcorner.
\]
For a theory $T$ with an $\mathcal E_a$-semidecidable axiom set,
fix a recursive Hilbert calculus with modus ponens and generalization,
and let $\mathrm{Prf}_T$ be its finite derivation trees whose conclusions
are sentences. Membership is
$\mathcal E_a$-semidecidable: syntax and inference checks are recursive,
and the finitely many nonlogical axiom tests can be run in parallel.
The validity of a derivation is independent of the name. Thus
$\mathrm{Prf}_T$ is an object of the splitting, and reading the root
conclusion gives a total arrow
\[
 c_T:\mathrm{Prf}_T\to S.
\]
Proposition~\ref{prop:first-order-proof-initial} characterizes this object
as the pullback of the initial derivation algebra to the sentence object.
Its range $t_T=\ran(c_T)$ is the partial identity on the theorem set.
The inverse image $\nu^*t_T$ is the refutation domain.

For $n\in L_a$, the sentence $c_n<1/2$ is an axiom of $T_a$.
For $n\in R_a$, the axiom $1/2<c_n$ and the ordered-field axioms
prove its negation. Consequently
\[
 1_{L_a}\le\theta^*t_{T_a},\qquad
 1_{R_a}\le\theta^*\nu^*t_{T_a}.
\]
By Corollary~\ref{cor:miller-cuts-inseparable}, these two source domains
are inseparable; the full inclusion into the restriction splitting
preserves this property. Finite domain joins and their inverse
images in the splitting are the corresponding set operations.
For a consistent extension $T\supseteq T_a$, we have
$t_{T_a}\le t_T$ and $t_T\nu^*t_T=0_S$.
Corollary~\ref{cor:domain-incompleteness} therefore makes
$t_T,\nu^*t_T$ inseparable and gives
$t_T\vee\nu^*t_T\ne1_S$. This is incompleteness of $T$.
\end{proof}

\section{The syntactic recursion category}\label{sec:syntactic}

Let $T$ be a consistent classical theory in the language of arithmetic
containing $\PA$. Write $\N$ for the natural numbers and $\omega$
for their formal numerical type. The other formal numerical types are
generated from $\omega$ by binary products and tagged coproducts.
Fix recursive numerical codings
whose bijection and structural equations are provable in $\PA$.
Fix a deterministic programming system whose functionality, composition,
and parameterization identities are provable in $\PA$.
For an index $e$ of a program from $X$ to $Y$, let $G_e(x,y)$ be its
$\Sigma_1$ output formula, defined by existential quantification over
a primitive-recursive computation-trace predicate.

The category $S_T$ has these formal types as objects and typed indices
modulo the equivalence relation
\begin{equation}\label{eq:uniform}
 e\equiv_T f
 \quad\Longleftrightarrow\quad
 T\vdash \forall x\forall y\,
       \bigl(G_e(x,y)\leftrightarrow G_f(x,y)\bigr).
\end{equation}
Write $[e]_T$ for the equivalence class of the typed index $e$.
This is the unprimed syntactic construction of Montagna and
Di Paola--Montagna \cite{Montagna,DPM}; Montagna's $S$ is $S_{\PA}$.

For a program operation $\operatorname{comp}$ implementing composition,
$\PA$ proves
\begin{equation}\label{eq:composition}
 G_{\operatorname{comp}(f,e)}(x,z)
 \leftrightarrow
 \exists y\,\bigl(G_e(x,y)\land G_f(y,z)\bigr).
\end{equation}
Substitution of uniformly equivalent graphs under this existential proves
that composition descends to \eqref{eq:uniform}. The identity and
associativity equations follow from the same graph identities.

The pointwise relation defining the primed systems is different:
two programs are identified when $T$ proves their partial computations
equal separately at each standard input. Montagna's $S'$ is the case
$T=\PA$ of Di Paola and Montagna's $S'_T$ \cite{Montagna,DPM}.
\emph{Pointwise provable equality and the failure of composition}
\cite[Corollary 2]{Composition} proves that the proposed composition
is not independent of representatives for any consistent c.e.\
extension $T$ of $\PA$. Thus these primed systems do not define
categories with that composition.

For a $\Sigma_1$ formula $C(x)$ on the numerical type $X$, let
$[C]_T:X\rightharpoonup X$ be the partial identity represented by
the program that searches for a witness to $C(x)$ and then returns $x$.
For two such formulas,
\begin{equation}\label{eq:domain-order}
 [C]_T\le[D]_T
 \quad\Longleftrightarrow\quad
 T\vdash\forall x\,(C(x)\longrightarrow D(x)).
\end{equation}
Conjunction and disjunction give meets and joins of restriction
idempotents. The restriction $\overline{[e]_T}$ and range $\ran([e]_T)$
are represented respectively by
\[
 \exists y\,G_e(x,y),\qquad
 R_e(y):=\exists x\,G_e(x,y).
\]
For the range property, a restriction idempotent $[C]_T$ receives
$[e]_T$ precisely when
$T$ proves $G_e(x,y)\to C(y)$ universally, equivalently when
$T$ proves $R_e(y)\to C(y)$ universally.
If $g,h:Y\rightharpoonup Z$ satisfy $g[e]_T=h[e]_T$, then $T$
proves equality of the composite graphs. A witness to $R_e(y)$
therefore gives equality of the graphs of representatives of $g$ and $h$
at $y$. Existential elimination in $T$ gives
$g\ran([e]_T)=h\ran([e]_T)$.
The same range programs represent ranges over every extension of $T$.

The product and tagged-coproduct programs supply the restriction
products and coproducts. The universal program defines
$\tau:\omega\times\omega\rightharpoonup\omega$.
For every $f:X\times\omega\rightharpoonup\omega$, formalized
parameterization gives an arrow $r:X\to\omega$ represented by a
$\PA$-provably total program, with $f=\tau(r\times1_\omega)$.
Thus $\tau$ is a Turing morphism. The chosen coding
isomorphisms make $S_T$ isotypical, and product with a zero arrow is zero.
Thus $S_T$ is a prodominical recursion category.

\section{Interpretation in models}\label{sec:models}

In $\Pfn$, restriction products are Cartesian products and coproducts
are tagged disjoint unions.

Fix a model $M\models T$. Interpret numerical types recursively by
\[
 X_M=|M|\quad(X=\omega),\qquad
 (X\times Y)_M=X_M\times Y_M,\qquad
 (X+Y)_M=X_M\amalg Y_M.
\]
The fixed arithmetic codings identify these sets with the corresponding
codes in $M$. A standard program index in a graph formula is interpreted
as its numeral in $M$. Define
\begin{equation}\label{eq:model-functor}
 F_M(X)=X_M,\qquad
 F_M([e]_T)(x)=y
 \quad\Longleftrightarrow\quad M\models G_e(x,y).
\end{equation}
All quantified traces and intermediate values in this formula range over
the whole model, including its nonstandard elements.

\begin{proposition}\label{prop:model}
For every $M\models T$, \eqref{eq:model-functor} defines a covariant
functor $F_M:S_T\to\Pfn$ preserving restrictions, zeros, binary
restriction products and coproducts, ranges, and finite joins of
restriction idempotents.
\end{proposition}
\begin{proof}
Functionality is provable in $\PA$, so the interpreted graph is a
partial function. Equation \eqref{eq:uniform} makes the assignment
independent of the representative. Equation \eqref{eq:composition},
interpreted in $M$, is exactly relational composition of the two partial
functions; their intermediate values range over $Y_M$. The identity
graph gives the identity function.

The definedness predicate $\exists y\,G_e(x,y)$ gives the interpreted
function's domain. Similarly,
\[
 \{y\in Y_M:M\models R_e(y)\}
 =\operatorname{im}(F_M[e]_T),
\]
so ranges are preserved. False predicates, conjunction, and disjunction
give zeros, intersections, and unions. The PA-provable product and
tagged-coproduct graph equations give the corresponding operations on
partial functions.
\end{proof}

\section{A faithful family and theory extension}\label{sec:family}

Let $\mathcal M_T$ be the set of all structures for the language of
arithmetic with underlying set $\N$ that satisfy $T$. Their arithmetic
operations need not be the standard ones, and the interpretation of a
numeral $\bar n$ need not be the carrier element $n$. Completeness and
the downward L\"owenheim--Skolem theorem imply
$\mathcal M_T\ne\varnothing$: every model of $T$ is infinite, and a
countably infinite model can be transported to the carrier $\N$.

Define the product category
\[
 \mathcal D_T=\prod_{M\in\mathcal M_T}\Pfn.
\]
An object is a family of sets $(A_M)_M$; an arrow is a family of partial
functions with the indicated sources and targets. Composition and the
restriction, range, restriction-product, and coproduct operations are coordinatewise.
The functors of Proposition~\ref{prop:model} combine to give
\[
 F_T:S_T\longrightarrow\mathcal D_T,\qquad
 F_T(X)=(X_M)_M,\quad F_T([e]_T)=(F_M[e]_T)_M.
\]

\begin{theorem}\label{thm:faithful}
The functor $F_T$ is faithful and preserves restrictions, zeros, binary
restriction products and coproducts, ranges, and finite joins of
restriction idempotents. The category
$\mathcal D_T$ is split and locally connected, satisfies weak choice,
and has inverse images preserving binary joins of restriction
idempotents. Moreover,
$F_T(1_\omega)\ne0$.
\end{theorem}
\begin{proof}
Preservation follows coordinatewise from Proposition~\ref{prop:model}.
If $[e]_T\ne[f]_T$, let $E$ denote the uniform graph-equality sentence
in \eqref{eq:uniform}. The theory $T\cup\{\neg E\}$ is consistent.
Completeness and downward L\"owenheim--Skolem give
$M\in\mathcal M_T$ in which the two graphs differ. Hence
$F_M([e]_T)\ne F_M([f]_T)$, so $F_T([e]_T)\ne F_T([f]_T)$.

A restriction idempotent in $\mathcal D_T$ is a family of partial
identities on subsets. Its splitting uses those subsets as the
coordinate sets. For an object $A=(A_M)_M$, take the partial identity
on a singleton $\{x\}\subseteq A_M$ at one coordinate $M$ and the zero
map elsewhere. This restriction idempotent is connected. These
restriction idempotents detect the definedness and value of every coordinate of a partial arrow. Hence
$\mathcal D_T$ is locally connected.

A partial monomorphism has injective coordinate functions: test its
defining property on singleton objects supported at one coordinate.
The partial inverses of these injections form a section, proving weak
choice. Inverse images commute with unions coordinatewise. Finally,
$\mathcal M_T$ is nonempty and every interpreted numerical object is
nonempty, so its identity differs from zero.
\end{proof}

For consistent computably enumerable $T\supseteq\PA$, the source $S_T$
is not locally connected \cite[Corollary 4.4, pp.~657--658]{DPM}.
The connected restriction idempotents used in the preceding proof
belong to $\mathcal D_T$ and need not be images of source restriction
idempotents.

For an inclusion $T\subseteq T'$ of consistent theories, define
\[
 Q_{T,T'}:S_T\to S_{T'},\qquad
 Q_{T,T'}X=X,\quad Q_{T,T'}[e]_T=[e]_{T'}.
\]
Provability is monotone, so this is a well-defined covariant functor.
It preserves restrictions, zeros, ranges, binary restriction products
and coproducts, and the chosen Turing morphism.
Since $\mathcal M_{T'}\subseteq\mathcal M_T$, restriction of families
defines a covariant projection
$\pi_{T,T'}:\mathcal D_T\to\mathcal D_{T'}$. Fixing type interpretations
once for each arithmetic structure gives a strictly commuting square:
\[
\begin{tikzcd}[column sep=large]
 S_T \arrow[r,"Q_{T,T'}"] \arrow[d,"F_T"'] &
 S_{T'} \arrow[d,"F_{T'}"]\\
 \mathcal D_T \arrow[r,"\pi_{T,T'}"'] & \mathcal D_{T'}.
\end{tikzcd}
\]
Both paths send a program to the family of its interpretations in the
models of $T'$. The identity and composition laws for the projections
and quotient functors follow from restriction of indices and the
unchanged program representatives.

\subsection{Lifting in the syntactic category}

\begin{lemma}\label{lem:switch}
Let $c,c'$ be complementary restriction idempotents on $X$ in $S_T$.
There is a total arrow $s:X\to X+X$ with
\[
 sc=i_1c,\qquad sc'=i_0c'.
\]
The arrow $\chi_{X,X}^{-1}F_Ts$ is the unique section of
$\nabla_{F_TX}(F_Tc'+F_Tc)$, where
$\chi_{X,X}:F_TX+F_TX\to F_T(X+X)$ is the coproduct isomorphism.
\end{lemma}
\begin{proof}
Choose $\Sigma_1$ representatives $C(x)=\exists u\,R(x,u)$ and
$C'(x)=\exists v\,R'(x,v)$ with primitive-recursive witness tests.
The theory $T$ proves universally their disjointness and coverage.
For $j=0,1$, write $i_jx$ for the fixed numerical code of the
tagged input, so that this coding represents the coproduct injection
$i_j:X\to X+X$.
Search increasing finite bounds for a witness to either predicate,
returning the tagged input $i_1x$ on finding a $C$-witness and $i_0x$
on finding a $C'$-witness. Fix a priority for simultaneous discoveries.
Let $e_s$ be an index of this program. The finite-search verification
in $\PA$ proves that the program halts
exactly when $C(x)\lor C'(x)$. Under the disjointness assumption it also
proves the exact graph equation
\[
 G_{e_s}(x,z)\leftrightarrow
 (C(x)\land z=i_1x)\lor(C'(x)\land z=i_0x).
\]
Thus $T$ proves this equation and universal halting. The equations for
the restrictions follow, giving the total arrow $s=[e_s]_T$ of $S_T$.

In each $M\models T$, the interpretations of $C,C'$ partition $X_M$.
The interpreted program sends each input to its specified tagged copy.
It is exactly the inverse of the partial injection
$\nabla_{X_M}(F_M(c')+F_M(c))$. Its graph is therefore the unique required
section.
\end{proof}

\begin{corollary}\label{cor:pair}
For every consistent $T\supseteq\PA$, the functor $F_T$ satisfies the
lifting hypothesis of Theorem~\ref{thm:local-lifting}.
Fix $U=\omega$, a Turing morphism $\tau:U\times U\rightharpoonup U$,
and a PA-provable coding isomorphism $b:U\to U+U$. Put
\[
 d=\tau\Delta_U,\qquad
 a_0=(bd)^*(1_U+0_U),\qquad a_1=(bd)^*(0_U+1_U).
\]
Then $a_0,a_1$ are inseparable in $S_T$. For every consistent
extension $T'\supseteq T$, their images under $Q_{T,T'}$ are the
identically defined inseparable pair in $S_{T'}$.
\end{corollary}
\begin{proof}
Theorem~\ref{thm:faithful} and Lemma~\ref{lem:switch} verify all the
hypotheses of Theorem~\ref{thm:local-lifting}. The quotient functor keeps
every program in the definition fixed, and the same argument applies
to consistent $T'$.
\end{proof}

The image $F_T(\tau)$ is not universal in $\mathcal D_T$: at each
model coordinate, countably many program parameters cannot enumerate
all partial functions on its countable numerical carrier. The target
restriction lattices are Boolean, so every image of a disjoint source
pair has a separator. Lemma~\ref{lem:switch} lifts the sections
associated with complementary source restriction idempotents;
an arbitrary target separator need not be the image of such an
idempotent.

\section{Essential incompleteness of arithmetic}\label{sec:arithmetic}

Let $\mathrm{Rec}$ be the classical recursion category on the formal
numerical types, with partial recursive functions as arrows and extensional
equality. Its restriction idempotents are partial identities on
computably enumerable subsets; the complemented ones are partial
identities on recursive subsets.

Use the universal program of Section~\ref{sec:syntactic}, whose
computation and parameterization identities are provable in $\PA$.
It defines a Turing morphism
$\tau:\omega\times\omega\rightharpoonup\omega$ in $\mathrm{Rec}$.
Put $d=\tau\Delta_{\omega}$ and choose the coproduct isomorphism
$\zeta:\omega\to\omega+\omega$ whose numerical function is
$\zeta(2m)=i_0m$, $\zeta(2m+1)=i_1m$.
Let $\alpha(x)$ and $\beta(x)$ be the $\Sigma_1$ formulas asserting
that $d(x)$ halts with an even or an odd output, respectively.
Their standard truth sets
\[
 A=\{n:\N\models\alpha(\bar n)\},\qquad
 B=\{n:\N\models\beta(\bar n)\}
\]
are the classical even-output and odd-output sets for self-application.
Their recursive inseparability is another instance of
Theorem~\ref{thm:local-lifting}. To see this, let
$J:\mathrm{Rec}\to\Pfn$ send a formal numerical type to its underlying
set and a partial recursive arrow to the same partial function.
It is a covariant range functor. The category $\Pfn$ satisfies
the target hypotheses: partial identities on singletons detect partial
functions, restriction idempotents split through their subsets of
definition, partial injections have partial inverses, and inverse images
preserve unions. Complementary restriction idempotents on a numerical type $W$
in $\mathrm{Rec}$ are partial identities on complementary recursive
subsets $C,C'$ of its carrier $J(W)$. The arrow $s:W\to W+W$
given numerically by $s(x)=i_1x$ for $x\in C$ and $s(x)=i_0x$
for $x\in C'$ is total recursive. Its image under $J$ is the
required section. This verifies the lifting hypothesis of
Theorem~\ref{thm:local-lifting} for $J$.
The theorem, with the chosen universal program and parity
isomorphism $\zeta$, gives precisely the pair $A,B$.
Each finite computation gives a PA-proof of its numeral instance, and $\PA$ proves
$\forall x\,\neg(\alpha(x)\land\beta(x))$.

Work in
$\operatorname{Split}(\mathrm{Rec})$, whose objects can be represented
by c.e.\ numerical carriers and whose arrows are partial recursive
functions with c.e.\ domains in those carriers. A restriction idempotent
is a partial identity on a c.e.\ subset. Joins of these idempotents
correspond to unions of the subsets, and inverse image along a total
arrow preserves these joins. Write $a,b$ for the partial
identities on the sets $A,B$ obtained above. They remain inseparable
in this splitting.

Let $S$ be the recursive set of codes of closed arithmetic sentences.
Negation and numeral substitution give total maps
\[
 \nu:S\to S,\qquad \nu(\ulcorner\varphi\urcorner)
     =\ulcorner\neg\varphi\urcorner,
 \qquad
 \theta:\N\to S,\qquad \theta(n)=\ulcorner\alpha(\bar n)\urcorner.
\]
For a c.e.\ theory $T\supseteq\PA$, fix a recursive Hilbert calculus
with modus ponens and generalization. Let $\mathrm{Prf}_T$ be the
carrier of finite derivation trees from $T$ whose conclusions are sentences.
This carrier is c.e.: the logical rules and formula syntax are
recursive, and each candidate proof has only finitely many axiom
membership tests in the c.e.\ axiom set of $T$. Reading the root
conclusion gives a total arrow
$c_T:\mathrm{Prf}_T\to S$ in the splitting. Its range
$t_T=\ran(c_T)$ is the partial identity on the c.e.\ set of theorems
of $T$.
Proposition~\ref{prop:first-order-proof-initial} gives the initial-algebra
and pullback characterization of $c_T$.

The finite computation proofs of $\alpha(\bar n)$ for $n\in A$ give
$a\le\theta^*t_{\PA}$. For $n\in B$, a finite computation proves
$\beta(\bar n)$ in $\PA$; the PA-provable disjointness of $\alpha$
and $\beta$ then proves $\neg\alpha(\bar n)$. Thus
\begin{equation}\label{eq:arithmetic-domain-inclusions}
 a\le\theta^*t_{\PA},\qquad
 b\le\theta^*\nu^*t_{\PA},\qquad
 t_{\PA}\le t_T.
\end{equation}

\begin{corollary}[Essential incompleteness of arithmetic]\label{cor:arithmetic}
The theorem and refutation sets of every consistent computably
enumerable extension of $\PA$ are recursively inseparable.
Consequently every such extension is incomplete.
\end{corollary}
\begin{proof}
Let $T\supseteq\PA$ be consistent and c.e.\ Consistency says that
no sentence and its negation are both theorems, so
$t_T\,\nu^*t_T=0_S$. Completeness would say that every sentence or
its negation is a theorem, giving $t_T\vee\nu^*t_T=1_S$.
The pair $a,b$ supplied by Theorem~\ref{thm:local-lifting}, together
with \eqref{eq:arithmetic-domain-inclusions}, satisfies the hypotheses
of Corollary~\ref{cor:domain-incompleteness}. That corollary makes
$t_T,\nu^*t_T$ inseparable and rules out the completeness equation.
A complemented restriction idempotent on the recursive sentence carrier
$S$ is a partial identity on a recursive subset of $S$, so this gives
recursive inseparability of the theorem and refutation sets.
\end{proof}

\subsection{Numeral instances and theory extension}

For a consistent c.e.\ extension $T$ of $\PA$, define
\[
 P_T=\{n:T\vdash\alpha(\bar n)\},\qquad
 N_T=\{n:T\vdash\neg\alpha(\bar n)\}.
\]
These are the carriers of $\theta^*t_T$ and $\theta^*\nu^*t_T$.
They are computably enumerable disjoint sets with $A\subseteq P_T$
and $B\subseteq N_T$, and they enlarge under theory extension.
Let $\operatorname{Pair}(\mathrm{Rec})$ have objects $(X,A,B)$ consisting
of a numerical type and two disjoint computably enumerable subsets.
An arrow $(X,A,B)\to(Y,P,N)$ is a total recursive function $f:X\to Y$
with $f(A)\subseteq P$ and $f(B)\subseteq N$. Identities and composition
are inherited from $\mathrm{Rec}$. The assignment
$T\mapsto(\omega,P_T,N_T)$, with the identity function for theory
inclusions, is a covariant functor from the poset of consistent
computably enumerable extensions of $\PA$ to
$\operatorname{Pair}(\mathrm{Rec})$.
The inclusions $A\subseteq P_T$ and $B\subseteq N_T$ give the morphism
$1_\omega:(\omega,A,B)\to(\omega,P_T,N_T)$ in this category.
A recursive separator pulls back along a total recursive function;
consequently $P_T,N_T$ are recursively inseparable.


For a restriction idempotent
$c=[C]_T$ on $\omega$, define the subset of the standard natural numbers
\begin{equation}\label{eq:observation}
 E_T(c)=\{n\in\N:T\vdash C(\bar n)\}.
\end{equation}
It is computably enumerable by proof enumeration and is independent of
the representative. The operation is monotone, preserves finite meets,
and preserves top and bottom. Let $\kappa_n:\omega\to\omega$ be the
total constant program with value $n$. Then
\begin{equation}\label{eq:faithful-observation}
 n\in E_T(c)
 \quad\Longleftrightarrow\quad c\kappa_n=\kappa_n
 \quad\Longleftrightarrow\quad
 F_T(c)F_T(\kappa_n)=F_T(\kappa_n).
\end{equation}
The first equivalence follows from the partial-identity graph equation;
the second follows from faithfulness of $F_T$.


The sets $E_T(c)$ are computably enumerable because the axioms of $T$
are computably enumerable. For $T=\operatorname{Th}(\N)$, the functor
$F_T$ still exists and $S_T$ is the classical category of partial
recursive functions on the formal numerical types, but $T$ is complete.
Thus inseparability in $S_T$ alone does not imply incompleteness of $T$.

Bernardi's earlier method uses precomplete equivalence relations.
In the numerical case, precompleteness means that every partial
recursive function agrees, where defined and modulo the equivalence
relation, with a total recursive function. Under the hypotheses of
his categorical theorem, nonempty disjoint unions of equivalence
classes are inseparable. His examples include PA-provable program
equality \cite[pp.~125--128]{Bernardi}. The argument above instead
uses Theorem~\ref{thm:local-lifting} to obtain the numerical pair and
finite PA-proofs to place it in the sets $P_T,N_T$.

\section{A recursion category from a three-point space}\label{sec:cosheaves}

We construct a recursion category $\mathscr R_X$ whose restriction
splitting contains a consistent propositional theory with no complete
consistent extension represented by a restriction idempotent.
The construction uses cosheaves on a three-point space.

For a topological space $X$, let $\mathcal O(X)$ be the category
of open sets and inclusions. A \emph{set-valued cosheaf} is a
covariant functor $C:\mathcal O(X)\to\mathbf{Set}$ for which
every open cover $(O_\lambda)_\lambda$ of $O$ gives a coequalizer
\cite[Definition 2.6]{CurryCosheaves}
\[
 \coprod_{\lambda,\mu} C(O_\lambda\cap O_\mu)
 \rightrightarrows \coprod_\lambda C(O_\lambda)
 \longrightarrow C(O).
\]
The two parallel maps are induced by the inclusions into the two
members of the cover; the empty cover requires
$C(\varnothing)=\varnothing$.
The \emph{Sierpi\'nski space} has two points and exactly three open
sets: the empty set, one singleton, and the whole space.
Let $X=\{a,b,c\}$ have opens
$\varnothing,\{c\},\{a,c\},\{b,c\},X$.
The open subspaces $\{a,c\}$ and $\{b,c\}$ are Sierpi\'nski spaces
with common open singleton $\{c\}$. Thus $X$ is obtained by gluing
two copies of the Sierpi\'nski space along that singleton.
Its category of set-valued cosheaves is equivalent to
$\mathbf{Set}^{J}$, where $J$ has objects $0,1,2$ and nonidentity
morphisms $0\to1$ and $0\to2$.
The equivalence sends a cosheaf $C$ to the functor with values
$C(\{c\}),C(\{a,c\}),C(\{b,c\})$ at $0,1,2$, respectively,
and morphisms induced by the inclusions of these open sets.
The cosheaf value $C(X)$ is the pushout of
$C(\{a,c\})\leftarrow C(\{c\})\to C(\{b,c\})$
(Proposition~\ref{prop:fork-cosheaves}).
Let $\mathcal C_{\mathrm{rec}}$ have computably enumerable subsets
of finite powers of $\N$ as objects. An arrow $f:A\to B$ is a
function computed by a partial recursive program defined at every
element of $A$, with values there in $B$. Equality of arrows is
equality of their functions on $A$.
Put $\mathcal E=\mathcal C_{\mathrm{rec}}^J$.
An object is a covariant functor $Y:J\to\mathcal C_{\mathrm{rec}}$,
given by the objects $Y_j=Y(j)$, for $j=0,1,2$, and the
morphisms $r_i^Y=Y(0\to i):Y_0\to Y_i$, for $i=1,2$.
A morphism $f:Y\to Z$ is a natural transformation: its components
are arrows $f_j:Y_j\to Z_j$ of $\mathcal C_{\mathrm{rec}}$ satisfying
$r_i^Zf_0=f_i r_i^Y$.
A \emph{c.e.\ subfunctor} $D\subseteq Y$ consists of c.e.\ subsets
$D_j\subseteq Y_j$ such that $r_i^Y(D_0)\subseteq D_i$.
Forgetting the numerical presentations gives a functor
$\mathcal E\to\mathbf{Set}^J$; the cosheaf equivalence above assigns
a set-valued cosheaf to each such functor and a natural transformation
to each morphism of $\mathcal E$.
In the notation $\operatorname{Par}(\mathcal C,\operatorname{Mono})$,
$\operatorname{Mono}$ denotes all monomorphisms of the category
$\mathcal C$. Form the partial-map category
\[
 \mathcal A=\operatorname{Par}(\mathcal E,\operatorname{Mono}).
\]
Its arrows $Y\rightharpoonup Z$ are isomorphism classes of spans
$Y\xleftarrow{m}D\xrightarrow{f}Z$ with $m$ monic; composition
uses pullback. Monomorphisms are componentwise injections.
Thus a partial arrow is equivalently a computably enumerable
subfunctor $D\subseteq Y$ together with a natural transformation
$D\to Z$. Its
restriction is the partial identity on $D$, and its range is the
partial identity on the subfunctor of component images.
Proposition~\ref{prop:fork-structure} proves these assertions and
the structural properties used below.

Define an object $V:J\to\mathcal C_{\mathrm{rec}}$ of $\mathcal E$
by the following objects $V_j=V(j)$ and morphisms
$r_i^V=V(0\to i)$:
\begin{equation}\label{eq:universal-fork}
 \begin{gathered}
 V_0=\N^4,\qquad V_1=V_2=\N^2,\\
 r_1^V:V_0\to V_1,\qquad r_1^V(k,z,a,b)=(k,a),\\
 r_2^V:V_0\to V_2,\qquad r_2^V(k,z,a,b)=(k,b).
 \end{gathered}
\end{equation}
We regard $V$ also as an object of $\mathcal A$, whose objects
are those of $\mathcal E$. Let $V^\#$ be the object of
$\mathcal E$ whose component $(V^\#)_j$ consists of nonempty
finite words in $V_j$. Its morphism $(V^\#)(0\to i)$ applies
$r_i^V$ letter by letter, for $i=1,2$.
Choose isomorphisms in $\mathcal E$
$V\cong V\times V$, $V\cong V+V$, and $V\cong V^\#$.
Together with their inverses, these form a \emph{frame} at $V$
\cite[Section 3]{Heller}.
Let $\mathscr R_X$ denote the full subcategory of $\mathcal A$
on the objects obtained from $V$ by finitely many applications of
the chosen binary restriction products, binary coproducts, and
nonempty-word formation.
The isomorphisms, constructed in Lemma~\ref{lem:fork-frames}, make
$\mathscr R_X$ isotypical.

Theorem~\ref{thm:fork-turing} constructs a Turing morphism
$\tau:V\times V\rightharpoonup V$ in $\mathcal A$.
For every object $Y$ of $\mathcal A$ and every arrow
$F:Y\times V\rightharpoonup V$ in $\mathcal A$, there is a
total arrow $c:Y\to V$ such that $\tau(c\times1_V)=F$.
By fullness this property holds in $\mathscr R_X$ when
$Y$ belongs to $\mathscr R_X$. Together with the inherited
prodominical structure, it makes $\mathscr R_X$ a recursion category.

Every object $Y$ of $\mathcal E$ admits a total map $m:Y\to V$ in
$\mathcal A$ with a partial inverse $r:V\rightharpoonup Y$
such that $rm=1_Y$ and $mr=\bar r$
(Lemma~\ref{lem:fork-dominance}). In this sense $V$ is a
\emph{universal object} of $\mathcal A$.
These embeddings give the equivalence
$\operatorname{Split}(\mathscr R_X)\simeq\mathcal A$
proved in Proposition~\ref{prop:fork-recursion}.

\subsection{Local connectedness and lifting}

For an object $Y$ of $\mathcal E$ and an element $y\in Y_0$, take
the subfunctor with components $\{y\}$, $\{r_1^Yy\}$, and
$\{r_2^Yy\}$. Its partial identity is a connected restriction
idempotent in $\mathcal A$. The same holds for a subfunctor with
a singleton in component $1$ or $2$ and empty sets in the other
components. These restriction idempotents detect the definedness and value
of every component of a partial arrow, so $\mathcal A$ is locally
connected. Its splitness, weak choice, ranges, coproducts, and
componentwise inverse-image laws are proved in
Proposition~\ref{prop:fork-structure}.
The full inclusion $\mathscr R_X\hookrightarrow\mathcal A$ is a
range functor.
For each complementary pair $c,c'$ on $V$, the section
$s_c:V\to V+V$ of $\nabla_V(c'+c)$ belongs to $\mathscr R_X$
by fullness. Thus this inclusion satisfies the lifting hypothesis of
Theorem~\ref{thm:local-lifting}, which gives an inseparable pair in
$\mathscr R_X$.
The category $\mathscr R_X$ is not restriction-equivalent to the
classical recursion category. Indeed, the terminal object of
$\mathcal E$ has five subobjects, whereas the terminal object of
$\mathcal C_{\mathrm{rec}}$ has two.
A restriction equivalence with the classical recursion category
would extend to the restriction splittings and then restrict to an
equivalence $\mathcal E\simeq\mathcal C_{\mathrm{rec}}$, contradicting
this invariant. Proposition~\ref{prop:fork-distinction} gives the
subobject calculation.

\subsection{Evaluation functors}

Let
$\mathcal K=\operatorname{Par}(\mathcal C_{\mathrm{rec}},\operatorname{Mono})$.
This is the restriction splitting of the classical recursion category:
objects are c.e.\ numerical carriers and arrows are partial recursive
functions with c.e.\ domains contained in their sources.
For $j\in\{0,1,2\}$, component evaluation defines
\begin{equation}\label{eq:endpoint-functor}
 E_j:\mathcal A\longrightarrow\mathcal K,\qquad
 E_j(Y)=Y_j,\qquad
 E_j(Y\xleftarrow{m}D\xrightarrow{f}Z)
   =(Y_j\xleftarrow{m_j}D_j\xrightarrow{f_j}Z_j).
\end{equation}
For the Turing morphism $\tau:V\times V\rightharpoonup V$ of
Theorem~\ref{thm:fork-turing}, write
$\tau_j=E_j(\tau):V_j\times V_j\rightharpoonup V_j$.

\begin{proposition}\label{prop:endpoint}
The covariant functors $E_j$ preserve restrictions, zeros, ranges,
binary restriction products and coproducts, and finite joins of
restriction idempotents.
For $i=1,2$, the Turing morphism
$\tau:V\times V\rightharpoonup V$ constructed in
Theorem~\ref{thm:fork-turing} is sent
to a Turing morphism on $\N^2$ in $\mathcal K$.
\end{proposition}
\begin{proof}
Component evaluation preserves identities, monomorphisms, and
pullbacks in the total categories. Consequently \eqref{eq:endpoint-functor}
is independent of the representative span and preserves partial-map
composition. The stated operations are represented componentwise by
partial identities, empty sets, images, products, tagged unions, and
finite unions, respectively.

Use the recursive bijection $\operatorname{encode}_2:\N^2\to\N$,
its inverse $\operatorname{decode}_2:\N\to\N^2$, and the acceptable numbering
$(\varphi_p)_{p\in\N}$ chosen in the construction of $\tau$ in
Appendix~\ref{app:fork-structure}. For $i=1,2$, its component has
the formula
\[
 (E_i\tau)((k,p),(\ell,x))
 =\operatorname{decode}_2\varphi_p(\operatorname{encode}_2(\ell,x)).
\]
For a partial recursive map
$F:Z\times\N^2\rightharpoonup\N^2$ with c.e.\ carrier $Z$,
the parameter theorem gives a total recursive map $s:Z\to\N$ such that
\[
 \varphi_{s(z)}
 =\operatorname{encode}_2\circ F(z,-)\circ\operatorname{decode}_2
 :\N\rightharpoonup\N
 \qquad(z\in Z).
\]
Thus $z\mapsto(0,s(z)):Z\to\N^2$ is the required total parameter map.
\end{proof}

For $i\in\{1,2\}$, let $H_i:J\to\mathcal C_{\mathrm{rec}}$
have component $\{*\}$ at $i$ and empty components elsewhere.
For $Y\in\mathcal E$ and $y\in Y_i$, let $p_y:H_i\to Y$ be the total morphism
whose component at $i$ sends $*$ to $y$. Every total morphism
$H_i\to Y$ has this form for a unique $y\in Y_i$. A partial
arrow from $H_i$ is either total or zero, since $H_i$ has just
two restriction idempotents. Hence postcomposition defines a partial
function on these total morphisms, defining
$\Phi_i:\mathcal A\to\Pfn$ by
\[
 \Phi_i(Y)=\operatorname{Tot}(\mathcal A)(H_i,Y),\qquad
 \Phi_i(f)(p)=fp\quad\hbox{when }fp\hbox{ is total}.
\]
If $fp$ is not total, it is zero, and every further composite
remains zero. If it is total, further postcomposition has exactly
the domain specified by the next partial function. This proves
$\Phi_i(gf)=\Phi_i(g)\Phi_i(f)$; identities are immediate.
Let $U:\mathcal K\to\Pfn$ send a c.e.\ numerical carrier to its
underlying set and a partial recursive arrow to the same partial
function. The bijections $p_y\mapsto y$ give a natural isomorphism
$\Phi_i\cong U E_i$.
Indeed, $fp_y=p_{E_i(f)(y)}$ when $E_i(f)(y)$ is defined, and $fp_y=0$
otherwise.
For a restriction idempotent $e$ on $Y$, put
\begin{equation}\label{eq:probe-domain}
 |e|_i:=\{y\in Y_i:ep_y=p_y\}
       =\{y\in Y_i:E_i(e)(y)=y\}.
\end{equation}
This is the c.e.\ subset on which $E_i(e)$ is the partial identity.

The functor $E_0$ does not preserve this Turing morphism's universality.
The total function $F:V_0\to V_0$ given by
$F(k,z,a,b)=(0,0,z,0)$ cannot equal $\tau_0(p,-)$ for any $p\in V_0$.
Inputs differing only in $z$ have the same image $(k,a)$ under
$r_1^V$. Naturality would therefore make their outputs have the
same image under $r_1^V$, whereas
$r_1^V F(k,z,a,b)=(0,z)$ depends on $z$.

\subsection{Propositional theories and component theorem sets}

Let $\zeta:V\to V+V$ be the coproduct isomorphism of
Lemma~\ref{lem:fork-frames}: it selects the summand by the parity
of the first coordinate. Define complementary restriction idempotents
$q_0,q_1:V\rightharpoonup V$ by
\[
 q_0=\zeta^{-1}(1_V+0_V)\zeta,\qquad
 q_1=\zeta^{-1}(0_V+1_V)\zeta.
\]
Put $d=\tau\Delta_V:V\rightharpoonup V$ and define
$a,b:V\rightharpoonup V$ by
\[
 a=d^*q_0,\qquad b=d^*q_1.
\]
These are the inseparable restriction idempotents supplied by
Theorem~\ref{thm:local-lifting} for the full inclusion
$\mathscr R_X\hookrightarrow\mathcal A$. They remain inseparable
in $\mathcal A$: fullness leaves the restriction idempotents on $V$, their order,
and their complements unchanged.
Let $A,B\subseteq V$ be the c.e.\ subfunctors on which $a,b$
are the partial identities, and write $A_j=A(j)$ and $B_j=B(j)$
for $j=0,1,2$. For $i\in\{1,2\}$, put $d_i=E_i(d)$.
By Proposition~\ref{prop:endpoint},
\begin{align*}
 A_i&=\{x:d_i(x)\text{ is defined with even first coordinate}\},\\
 B_i&=\{x:d_i(x)\text{ is defined with odd first coordinate}\}.
\end{align*}
These sets are recursively inseparable by the numerical instance
of Theorem~\ref{thm:local-lifting} used in
Section~\ref{sec:arithmetic}. Here the universal object is $\N^2$,
the Turing morphism is $E_i\tau$, and the coproduct isomorphism
is $E_i\zeta$, using the componentwise identification of
$E_i(V+V)$ with $\N^2+\N^2$. Proposition~\ref{prop:endpoint}
verifies universality, and the inclusion of partial recursive
functions into $\Pfn$ supplies the same range functor and lifts.

We construct a theory in classical propositional logic from the
disjoint restriction idempotents $a,b$. Define the polynomial functor
$P:\mathcal E\to\mathcal E$ by
$P(Z)=V+Z+(Z\times Z)$ and
$P(f)=1_V+f+(f\times f)$ for morphisms $f$ of $\mathcal E$.
Its initial algebra $\iota:P(S)\to S$ is constructed from finite
formula trees. Initiality means that every morphism
$\alpha:P(Z)\to Z$ of $\mathcal E$ admits a unique morphism $u:S\to Z$
of $\mathcal E$
with $u\iota=\alpha P(u)$.
The three summands give constructors for atoms, negation, and implication.
Write $\theta:V\to S$ for the total atom map
and $\nu:S\to S$ for the total negation constructor.
Let $\mathrm{Prf}$ be the object of finite proofs for the rules below
and $c:\mathrm{Prf}\to S$ its total conclusion map.
Proposition~\ref{prop:fork-proof-initial} characterizes $(\mathrm{Prf},c)$
as an initial algebra in $\mathcal E/S$.
The restriction idempotent $t_0=\ran(c)$ represents the theorems.
For $j=0,1,2$, write $S_j=S(j)$, and let $\theta_j$, $\nu_j$,
and $\Rightarrow_j$ be the components of the atom, negation,
and implication morphisms. At component $j$, the logical axiom
schemes, with $\alpha,\beta,\gamma\in S_j$, are
\[
\begin{gathered}
 \alpha\Rightarrow_j(\beta\Rightarrow_j\alpha),\\
 (\alpha\Rightarrow_j(\beta\Rightarrow_j\gamma))
   \Rightarrow_j((\alpha\Rightarrow_j\beta)\Rightarrow_j
                         (\alpha\Rightarrow_j\gamma)),\\
 (\nu_j\beta\Rightarrow_j\nu_j\alpha)
       \Rightarrow_j(\alpha\Rightarrow_j\beta).
\end{gathered}
\]
Modus ponens at component $j$ concludes $\beta$ from $\alpha$
and $\alpha\Rightarrow_j\beta$. The additional axioms are
$\theta_j(x)$ for $x\in A_j$ and $\nu_j\theta_j(y)$ for $y\in B_j$.
For a restriction idempotent $t$ on $S$, denote by $D_t\subseteq S$
the c.e.\ subfunctor on which $t$ is the partial identity.
Theorem~\ref{thm:fork-proof-presentation} constructs this syntax and
proof object, proves that $D_{t_0}$ is the least
subfunctor closed under modus ponens and containing the axioms, and proves
consistency. In particular, its one-node proofs give
\begin{equation}\label{eq:atomic-inequalities}
 a\le\theta^*t_0,\qquad
 b\le\theta^*\nu^*t_0,\qquad t_0\,\nu^*t_0=0.
\end{equation}
A \emph{theory extending $t_0$} is a restriction idempotent $t\ge t_0$
on $S$ in $\mathcal A$ such that $D_t$ is closed under modus ponens.
It is consistent when $t\,\nu^*t=0$, and complete when
$t\vee\nu^*t=1_S$. The set $(D_{\nu^*t})_j$ consists of the
formulas whose negations belong to $(D_t)_j$.

\begin{theorem}\label{thm:cosheaf-observation}
Fix $i\in\{1,2\}$. For every consistent theory $t$ extending
$t_0=\ran(c)$ on $S$ in $\mathcal A$, define
\[
 P_{i,t}=(E_i\theta)^{-1}(|t|_i),\qquad
 N_{i,t}=(E_i\theta)^{-1}(|\nu^*t|_i).
\]
The restriction idempotents $t,\nu^*t$ are inseparable in $\mathcal A$.
The sets $P_{i,t},N_{i,t}$ are disjoint c.e.\ subsets of $\N^2$ satisfying
$A_i\subseteq P_{i,t}$ and $B_i\subseteq N_{i,t}$.
They are recursively inseparable. The assignment
\[
 t\longmapsto(\omega\times\omega,P_{i,t},N_{i,t})
\]
with the identity of $\omega\times\omega$ for extension inclusions is a
covariant functor from the poset of consistent theories extending
$t_0$ to $\operatorname{Pair}(\mathrm{Rec})$.
Every such theory is incomplete.
\end{theorem}
\begin{proof}
The category $\mathcal A$ has restriction zeros and joins of restriction idempotents
preserved by total inverse image. Its restriction idempotents $a,b$ are inseparable,
and \eqref{eq:atomic-inequalities} gives the hypotheses of
Corollary~\ref{cor:domain-incompleteness}. Hence $t,\nu^*t$ are
inseparable, and $t$ is incomplete.

The atom map $E_i\theta:\N^2\to S_i$ is total recursive.
The component theorem and refutation sets $|t|_i$ and $|\nu^*t|_i$
are c.e.; their inverse images under the total recursive map
$E_i\theta$ are therefore c.e. Consistency and preservation of meets and
zero make them disjoint. Applying $E_i$ to
\eqref{eq:atomic-inequalities} and using $t_0\le t$ gives
the two inclusions. A recursive separator would consequently
separate $A_i,B_i$, which is impossible.

If $t\le t'$, monotonicity gives inclusions of both component
sets, so the identity of $\omega\times\omega$ is the required pair morphism.
The functor laws follow from the identity maps.
\end{proof}

\appendix

\section{The three-point cosheaf construction}\label{app:fork-structure}

\subsection{Cosheaves and functor categories}

Write $D=\{c\}$, $L=\{a,c\}$, and $R=\{b,c\}$ for the nonempty
proper opens of $X$. Morphisms of the set-valued cosheaves defined in
Section~\ref{sec:cosheaves} are natural transformations.

\begin{proposition}\label{prop:fork-cosheaves}
Restriction to $D,L,R$ defines an equivalence
$\operatorname{CoSh}(X,\mathbf{Set})\simeq\mathbf{Set}^{J}$,
where $J=(1\leftarrow0\to2)$. An object of $\mathbf{Set}^{J}$
consists of sets $C,A,B$ and functions $r:C\to A$, $s:C\to B$.
The inverse assigns to this functor the cosheaf whose values at
$D,L,R,X$ are respectively $C,A,B,A\amalg_C B$.
\end{proposition}
\begin{proof}
Every cover of $D$, $L$, or $R$ contains the covered open. A cover
of $X$ either contains $X$ or contains both $L$ and $R$. Thus the
only additional cosheaf condition is that the following square be
a pushout:
\[
\begin{tikzcd}
C\arrow[r,"s"]\arrow[d,"r"']&B\arrow[d]\\
A\arrow[r]&A\amalg_C B.
\end{tikzcd}
\]
Any other members of a cover factor through these maximal members,
so their identifications are already imposed by this pushout.
A natural transformation induces a unique map of pushouts.
This defines the inverse functor on morphisms; the same universal
property gives both natural isomorphisms required for an equivalence.
\end{proof}

For an object $Y$ of $\mathcal E$, the pushout has a
quotient presentation on the tagged carrier $Y_1+Y_2$: identify
$(1,r_1^Yy)$ with $(2,r_2^Yy)$ for every $y\in Y_0$ and take the
equivalence relation generated by these pairs. This relation is c.e.,
since finite paths in its c.e.\ edge relation can be enumerated.
A morphism of $\mathcal E$ acts on representatives and respects the
relation, and hence induces the component at $X$ of the associated
cosheaf morphism.

Products and subobjects are taken in $\mathcal E$. In
particular, the cosheaf associated with $Y\times Z$ has value at $X$
\[
 (Y_1\times Z_1)\amalg_{Y_0\times Z_0}(Y_2\times Z_2).
\]
This need not be the product of the values at $X$ of
the cosheaves associated with $Y$ and $Z$. Write $0$, $1$, and $2$
for empty, singleton, and two-element numerical carriers, respectively.
For example,
the object $H=(1\leftarrow0\to1)$ of $\mathcal E$ satisfies
$H\times H\cong H$,
and the cosheaf associated with $H$ has two elements at $X$,
whereas the Cartesian square of that set has four.
Likewise $H\to(1\leftarrow1\to1)$ is a monomorphism in $\mathcal E$,
although the associated cosheaf morphism has component $2\to1$ at $X$.
Thus the
subfunctors used as partial-map domains are categorical subobjects
of the cosheaves, not necessarily subsets at every open.

A \emph{regular epimorphism} is a coequalizer of a pair of arrows.
A category is \emph{regular} if it has finite limits, every arrow
factors as a regular epimorphism followed by a monomorphism, and
regular epimorphisms are preserved by pullback. The monic part of
this factorization is its \emph{regular image}. The \emph{kernel
pair} of $q:Y\to Z$ is the pair of projections
$Y\times_ZY\rightrightarrows Y$.
Write $\mathcal E/Y$ for the slice category of arrows with target $Y$,
whose morphisms are commuting triangles over $Y$.
Finite coproducts are \emph{extensive} when, for every finite family
$(Y_k)_k$, copairing defines an equivalence
\[
 \prod_k(\mathcal E/Y_k)\longrightarrow
 \mathcal E/\Bigl(\coprod_kY_k\Bigr).
\]

\begin{proposition}\label{prop:fork-total}
The categories $\mathcal C_{\mathrm{rec}}$ and
$\mathcal E=\mathcal C_{\mathrm{rec}}^J$ have finite limits and
finite coproducts. The category $\mathcal E$ is regular and
extensive. Its monomorphisms are componentwise injections, and its
regular images are the componentwise c.e.\ images.
\end{proposition}
\begin{proof}
The terminal numerical carrier is a singleton. Products use tuples,
and an equalizer of $f,g:Y\to Z$ has carrier
$\{y\in Y:f(y)=g(y)\}$. This carrier is c.e.: enumerate $Y$, run
both computations, and retain equal outputs. Its inclusion is
recursive. The universal maps into products and equalizers use the
given recursive maps and tuple formation, and are total on their
specified source carriers. This proves the finite-limit universal
properties in $\mathcal C_{\mathrm{rec}}$.

For a total arrow $f:Y\to Z$, enumerate $Y$ and its computed
values to enumerate $f(Y)$. The factorization
$Y\to f(Y)\hookrightarrow Z$ is therefore a factorization in
$\mathcal C_{\mathrm{rec}}$. Every total surjection
$q:Y\to Z$ has a partial recursive section total on $Z$: on input
$z$, enumerate $Y$ and search for a computed value $q(y)=z$.
If $u:Y\to W$ is constant on the fibers of $q$, composing $u$
with this section gives the unique arrow $v:Z\to W$ of
$\mathcal C_{\mathrm{rec}}$ with $vq=u$. Thus $q$ is the
coequalizer of its kernel pair.
The pullback of a surjection is surjective, since a preimage of
each element can be paired with that element. These pullbacks
have c.e.\ tuple carriers by the finite-limit construction.
This proves regularity and pullback stability of the image
factorizations in the numerical category.

Limits and tagged finite coproducts in $\mathcal E$ are formed
componentwise. Their universal maps are natural because their
components are uniquely determined by the corresponding numerical
universal properties. For a natural transformation $f:Y\to Z$,
the equations $r_i^Zf_0=f_i r_i^Y$ show that its component images
form a subfunctor of $Z$.
For a componentwise surjection $q:Y\to Z$, a natural transformation
$u:Y\to W$ coequalizing its kernel pair descends componentwise.
If $v_jq_j=u_j$ are the descended maps, then
\[
 r_i^Wv_0q_0=r_i^Wu_0=u_i r_i^Y
 =v_iq_i r_i^Y=v_i r_i^Zq_0.
\]
Surjectivity of $q_0$ gives $r_i^Wv_0=v_i r_i^Z$, so the
descent is a natural transformation. These image factorizations are stable
under pullback because the corresponding assertion holds in
each component. They give regularity of $\mathcal E$.

The three objects of $\mathcal E$
\[
 H_0=(1\leftarrow1\to1),\qquad
 H_1=(1\leftarrow0\to0),\qquad
 H_2=(0\leftarrow0\to1)
\]
detect component elements: total maps $H_j\to Y$ are in bijection
with elements of $Y_j$. These maps are recursive because every
nonempty component of their source $H_j$ is a singleton.
If a morphism of $\mathcal E$ is monic,
testing it on these maps proves that each component is injective.
The converse follows by cancelling those component injections.
An injective arrow of $\mathcal C_{\mathrm{rec}}$ has a partial recursive inverse
on its c.e.\ image, obtained by preimage search.

Finally, for a map $Y\to A+B$, the inverse images of the two
tags form c.e.\ subfunctors $Y^A,Y^B$ and give an isomorphism
$Y^A+Y^B\cong Y$. Maps over $A+B$ restrict uniquely to maps
over $A$ and $B$, and two such maps combine uniquely by tagged
copairing. This proves the slice equivalence expressing extensivity.
\end{proof}

\subsection{Partial maps, ranges, and connected restriction idempotents}

\begin{proposition}\label{prop:fork-structure}
The category $\mathcal A=\operatorname{Par}(\mathcal E,
\operatorname{Mono})$ is a split prodominical restriction category
with binary restriction products, restriction coproducts, weak
choice, and ranges satisfying the agreement law. The range of an arrow
represented by $Y\xleftarrow{m}D\xrightarrow{v}Z$ is the partial
identity on the regular image of $v$.
Joins of restriction idempotents are partial identities on finite
unions of their subfunctors, and inverse image along every partial
arrow preserves binary meets and finite joins. Ranges commute with products and
are stable under pullback along total arrows.
\end{proposition}
\begin{proof}
Replace the monic leg of a representing span by its image
inclusion, using its computable inverse on that image. A partial
arrow $f:Y\rightharpoonup Z$ is then represented by a c.e.\ subfunctor
$D_f\subseteq Y$ and a natural transformation $D_f\to Z$.
For $j\in\{0,1,2\}$, write $f_j:Y_j\rightharpoonup Z_j$ for the
partial recursive function with domain $(D_f)_j$ whose restriction
to $(D_f)_j$ is the $j$-component of this natural transformation.
Use the same convention for other partial arrows.
If $g:Z\rightharpoonup W$ has domain subfunctor $D_g$, the domain of $gf$
at component $j$ is
\[
 \{y\in(D_f)_j:f_j(y)\in(D_g)_j\},
\]
and its value there is $g_j(f_j(y))$. These carriers are c.e.\
and form a subfunctor by naturality. This is precisely the
pullback composition of the representing spans. Associativity
and identities follow from these domain and value formulas.

The restriction $\bar f$ is the partial identity on $D_f$.
The restriction axioms reduce componentwise to the corresponding
identities of partial functions. Every restriction idempotent
splits through its domain subfunctor: the inclusion is total,
and its partial inverse is the identity with that domain.
Empty domains give restriction zeros. The restriction product
uses componentwise Cartesian products; a pairing is defined on
the intersection of the two domains and has the ordered pair
of values there. Its projection identities and uniqueness are
those in the definition of a restriction product. Coproducts
use tagged carriers; partial copairing uses the tagged union
of the two domains. A product with an empty domain is empty,
which proves the prodominical zero law.

The component images $f_j((D_f)_j)$ form a c.e.\ subfunctor
$I_f\subseteq Z$. Its partial identity receives $f$. If a restriction
idempotent receives $f$, its subfunctor contains each component image,
and hence contains $I_f$.
This proves the range property and agreement with the regular
image of $D_f\to Z$. For $g,h:Z\rightharpoonup W$ with
$gf=hf$, choose a preimage of each $z\in(I_f)_j$. Equality of the two partial composites shows
that $g_j$ is defined at $z$ exactly when $h_j$ is, and their
values there agree. Hence $g\ran(f)=h\ran(f)$, including
equality of domains.

For a total $u:W\to Z$, the projection
$D_f\times_Z W\to W$ has component image
$u_j^{-1}(f_j((D_f)_j))$. Thus its range is $u^*\ran(f)$,
which is the asserted pullback law. The image of a product
of two representing morphisms is the product of their images, giving
$\ran(f\times g)=\ran(f)\times\ran(g)$.
Intersections and finite unions of domain subfunctors are again
c.e.\ subfunctors; their partial identities give the meets and joins
of the corresponding restriction idempotents. For a restriction idempotent
$e$ on $Z$ with domain subfunctor $D_e$, the inverse image $f^*e$
has component carrier
\[
 \{y\in(D_f)_j:f_j(y)\in(D_e)_j\}.
\]
This formula proves preservation of binary meets and finite joins,
including the empty join.

Suppose that $p:Y\rightharpoonup Z$ is a partial monomorphism
and let $D_p$ be its domain subfunctor. If two elements of $(D_p)_j$
have the same image, their corresponding total maps
$H_j\to D_p\to Y$ have equal composites with $p$.
Partial monicity makes the original maps equal, since their
images lie in $D_p$. Thus each $p_j$ is injective on $(D_p)_j$.
Preimage search gives recursive inverse components on $I_p$;
they form a natural transformation by injectivity and naturality of
the representing morphism $D_p\to Z$.
The resulting partial arrow $s:Z\rightharpoonup Y$ has domain
$I_p$ and satisfies $ps=\bar s$ and $psp=p$.
This proves weak choice.
\end{proof}

\begin{proposition}\label{prop:fork-local}
The category $\mathcal A$ is locally connected.
\end{proposition}
\begin{proof}
Let $Y$ be an object of $\mathcal E$. For $y\in Y_0$, take the
finite subfunctor $D\subseteq Y$ with
\[
 D_0=\{y\},\qquad D_1=\{r_1^Yy\},\qquad D_2=\{r_2^Yy\}.
\]
The partial identity on $D$ is connected: in any partition of $D$
into two disjoint subfunctors, the part containing $y$ in component
$0$ must contain its images in components $1$ and $2$, and hence
must equal $D$. The partial identity on a singleton subfunctor
in component $i\in\{1,2\}$, with the other components empty,
is connected as well.

Suppose that partial arrows $f,g:Y\rightharpoonup Z$ agree
after composition with every connected restriction idempotent.
Restriction to the
finite subfunctor generated by $y\in Y_0$ detects whether
$f_0$ and $g_0$ are defined at $y$, and compares their values
when defined. Partial identities on singleton subfunctors in
components $1$ and $2$ do the same for $f_1,g_1$ and $f_2,g_2$.
Thus all three component partial functions agree, proving $f=g$.
\end{proof}

\subsection{A universal object and a frame}

Let $V$ be the object of $\mathcal A$ defined in
\eqref{eq:universal-fork}.
A \emph{total restriction embedding} $m:Y\to V$ is a total
arrow with a partial inverse $m^\dagger:V\rightharpoonup Y$
satisfying $m^\dagger m=1_Y$ and $mm^\dagger=\overline{m^\dagger}$.

\begin{lemma}\label{lem:fork-dominance}
Every object of $\mathcal E$ has a total restriction embedding
into $V$.
\end{lemma}
\begin{proof}
Let $Y$ be an object of $\mathcal E$. For $j=0,1,2$, choose a
recursive injection $c_j:Y_j\to\N$
with a partial recursive inverse on its image, using a tuple
coding of the ambient finite power. Define
\[
 m_i(y_i)=(0,c_i(y_i))\quad(i=1,2),\qquad
 m_0(y)=(0,c_0(y),c_1(r_1^Yy),c_2(r_2^Yy)).
\]
These maps are total and injective, and satisfy
$r_i^Vm_0=m_i r_i^Y$ for $i=1,2$.
Their image is a c.e.\ subfunctor. The inverses on that image
are natural and give the required partial inverse.
\end{proof}

For an object $Y$ of $\mathcal E$, let $Y^\#$ have at component $j$ the
nonempty finite words over $Y_j$; the morphism $Y^\#(0\to i)$
applies $Y(0\to i)$ to each letter. Finite words over a c.e.\ carrier form
a c.e.\ carrier. A morphism of $\mathcal E$ acts on words letterwise.
A partial arrow in $\mathcal A$ acts on precisely the words all of whose
letters belong to its domain, also letterwise. These formulas
preserve identities and composition and define a covariant functor
$(-)^\#:\mathcal A\to\mathcal A$. Concatenation
$Y^\#\times Y^\#\to Y^\#$ and singleton insertion
$Y\to Y^\#$ are morphisms of $\mathcal E$. For an internal semigroup $M$
in $\mathcal E$ and a morphism $f:Y\to M$, multiplication of the
$f$-images of the letters gives the unique semigroup morphism
$\widehat f:Y^\#\to M$ whose composite with singleton insertion is $f$.

Fix a recursive pairing bijection
$\pi:\N^2\to\N$ and recursive tuple bijections
$\gamma_n:\N^n\to\N$, with their inverses computable uniformly
for $n\ge1$. Encode a nonempty word $(k_1,\ldots,k_n)$ by
$\pi(n-1,\gamma_n(k_1,\ldots,k_n))$.

\begin{lemma}\label{lem:fork-frames}
There are isomorphisms in $\mathcal E$
$V\cong V\times V$, $V\cong V+V$, and $V\cong V^\#$.
\end{lemma}
\begin{proof}
For the product isomorphism, decode each of $k,z,a,b$ into a
pair using $\pi^{-1}$. Send $(k,z,a,b)$ to the resulting pair
of four-tuples. At component $i\in\{1,2\}$, decode $(k,x)$ into
the resulting pair of two-tuples. The naturality equations hold
coordinatewise, and coordinatewise pairing gives the inverse.

For the coproduct isomorphism, write $k=2q+t$ with
$t\in\{0,1\}$. Send a tuple with first coordinate $k$ to
summand $t$, replace $k$ by $q$, and retain its remaining
coordinates. This rule is the same at all three components.
The inverse replaces $q$ in summand $t$ by $2q+t$.

For the word isomorphism, decode $k$ as the nonempty word
$(k_1,\ldots,k_n)$. Decode $z,a,b$ by $\gamma_n^{-1}$,
and send $(k,z,a,b)$ to
\[
 ((k_1,z_1,a_1,b_1),\ldots,(k_n,z_n,a_n,b_n)).
\]
At component $i\in\{1,2\}$, decode $k$ in the same way and
decode $x$ by $\gamma_n^{-1}$, obtaining $((k_j,x_j))_{j=1}^n$.
These maps form a natural transformation. Recovering the word of
first coordinates and applying the tuple encodings to the
remaining coordinate words gives total recursive inverses.
\end{proof}

\subsection{A Turing morphism}\label{app:fork-turing}

We construct a Turing morphism $\tau:V\times V\rightharpoonup V$
in $\mathcal A$ by giving a monomorphism $T\hookrightarrow V\times V$
and a natural transformation $T\to V$.

Fix recursive bijections
$\operatorname{encode}_n:\N^n\to\N$, with recursive inverses
$\operatorname{decode}_n$, for each positive integer $n$.
We use an acceptable numbering
$\varphi:\N\times\N\rightharpoonup\N$ of the partial recursive
functions and write $\varphi_p(t)=\varphi(p,t)$.
The specialization property of this numbering is the following:
if $m\geq0$, $n\geq1$, and
$G:\N^m\times\N^n\rightharpoonup\N$ is partial recursive,
there is a total recursive map $s:\N^m\to\N$ such that
\[
 \varphi_{s(y)}(\operatorname{encode}_n x)\simeq G(y,x).
\]
Here $\simeq$ means equality of partial values,
including simultaneous undefinedness. The case $m=0$ is a
constant program index. This is the usual $s$-$m$-$n$ theorem.

\begin{theorem}\label{thm:fork-turing}
The object $V$ of
$\mathcal A=\operatorname{Par}(\mathcal C_{\mathrm{rec}}^J,
\operatorname{Mono})$ admits a Turing morphism
$\tau:V\times V\rightharpoonup V$: for every object $Y$ of
$\mathcal A$ and every arrow $F:Y\times V\rightharpoonup V$
in $\mathcal A$, there is a total arrow
$c:Y\to V$ such that
\[
 \tau(c\times1_V)=F.
\]
At component $i\in\{1,2\}$, this Turing morphism is given by
\begin{equation}\label{eq:fork-turing-endpoint}
 \tau_i((k,p),(\ell,x))
 =\operatorname{decode}_2
     \varphi_p(\operatorname{encode}_2(\ell,x)).
\end{equation}
\end{theorem}

\begin{proof}
For $i=1,2$, let $T_i\subseteq V_i\times V_i$ be the c.e.\ set
on which the numerical computation in
\eqref{eq:fork-turing-endpoint} halts, and define
$\tau_i:V_i\times V_i\rightharpoonup V_i$ with domain $T_i$
by that equation. The first coordinate $k$
of the program argument has no effect on this computation.

At component $0$, write the program and input as
\[
 P=(k,p_0,p_1,p_2)\in V_0,
 \qquad x=(\ell,x_0,x_1,x_2)\in V_0.
\]
Define partial recursive maps
$C:V_0\times V_0\rightharpoonup V_0$ and
$Q_i:V_0\times V_0\rightharpoonup V_i$, for $i=1,2$,
by the following computations, run in parallel:
\begin{align*}
 C(P,x)&\simeq
   \operatorname{decode}_4
      \varphi_{p_0}(\operatorname{encode}_4 x),\\
 Q_i(P,x)&\simeq
   \operatorname{decode}_2
      \varphi_{p_i}(\operatorname{encode}_2(\ell,x_i))
      \qquad(i=1,2).
\end{align*}
Let $T_0\subseteq V_0\times V_0$ consist of those pairs $(P,x)$
for which all three computations halt and
\begin{equation}\label{eq:fork-turing-check}
 Q_i(P,x)=r_i^V C(P,x)\qquad(i=1,2).
\end{equation}
Define $\tau_0:V_0\times V_0\rightharpoonup V_0$ with domain
$T_0$ by $\tau_0(P,x)=C(P,x)$.
Dovetailing the three computations and checking equality of their
numerical outputs enumerates $T_0$ and computes $\tau_0$ on it.

If $(P,x)\in T_0$, then
\[
 (r_i^V P,r_i^V x)=((k,p_i),(\ell,x_i))\in T_i.
\]
Thus $T=(T_0,T_1,T_2)$ is a c.e.\ subfunctor of $V\times V$.
Equation~\eqref{eq:fork-turing-check} gives
\[
 r_i^V\tau_0(P,x)
   =\tau_i(r_i^V P,r_i^V x)
 \qquad((P,x)\in T_0).
\]
The restrictions $\tau_j|_{T_j}:T_j\to V_j$, for $j=0,1,2$,
therefore form a morphism $T\to V$ in $\mathcal E$.
Together with the inclusion $T\hookrightarrow V\times V$,
they define the required partial arrow $\tau$.

Now let $F:Y\times V\rightharpoonup V$ be any partial arrow.
Choose its representation by a c.e.\ subfunctor
$D\subseteq Y\times V$ and a morphism $D\to V$ in
$\mathcal E$. Its components are the restrictions
$F_j|_{D_j}:D_j\to V_j$, where $j=0,1,2$.
Put $d_0=4$ and $d_1=d_2=2$, so that $V_j=\N^{d_j}$.
Choose $m_j\geq0$ with $Y_j\subseteq\N^{m_j}$.
For each component define a partial recursive numerical function
\[
 G_j:\N^{m_j}\times\N^{d_j}\rightharpoonup\N
\]
as follows. On input $(y,x)$, wait until $(y,x)$ appears in an
enumeration of $D_j$; then compute $F_j(y,x)$ and return its
$\operatorname{encode}_{d_j}$-code. This function has domain
exactly $D_j$.
Specialization gives total recursive maps
$s_j:\N^{m_j}\to\N$ satisfying
\begin{equation}\label{eq:fork-specialization}
 \operatorname{decode}_{d_j}
   \varphi_{s_j(y)}(\operatorname{encode}_{d_j}x)
 \simeq F_j(y,x),
 \qquad\text{with domain }D_j.
\end{equation}

Define the component maps of $c:Y\to V$ by
\begin{align*}
 c_i(y_i)&=(0,s_i(y_i)) &&(i=1,2),\\
 c_0(y)&=(0,s_0(y),s_1(r_1^Yy),s_2(r_2^Yy)).
\end{align*}
These maps are recursive and total on their respective carriers:
the maps $s_j$ are total on the ambient numerical powers and the
maps $r_i^Y$ are total on $Y_0$. Moreover,
\[
 r_i^Vc_0(y)=(0,s_i(r_i^Yy))=c_i(r_i^Yy),
\]
so $c$ is a natural transformation in $\mathcal E$, hence total in
$\mathcal A$.

At component $i\in\{1,2\}$, equations
\eqref{eq:fork-turing-endpoint} and
\eqref{eq:fork-specialization} show that the partial function
\[
 (y_i,v_i)\longmapsto\tau_i(c_i(y_i),v_i)
 :Y_i\times V_i\rightharpoonup V_i
\]
has domain $D_i$ and agrees there with $F_i$.
At component $0$, fix $(y,x)\in Y_0\times V_0$.
If $(y,x)\in D_0$, the computation $C(c_0(y),x)$ returns $F_0(y,x)$.
Since $D$ is a subfunctor,
\[
 (r_i^Yy,r_i^Vx)\in D_i\qquad(i=1,2).
\]
The computations $Q_i(c_0(y),x)$, for $i=1,2$, consequently
halt, and naturality of the morphism $D\to V$ gives
\[
 F_i(r_i^Yy,r_i^Vx)=r_i^VF_0(y,x).
\]
Both tests in \eqref{eq:fork-turing-check} succeed, so
$\tau_0(c_0(y),x)=F_0(y,x)$.
Conversely, if $\tau_0(c_0(y),x)$ is defined, the computation
$C(c_0(y),x)$ has halted. Equation~\eqref{eq:fork-specialization}
then implies $(y,x)\in D_0$.
Thus $\tau(c\times1_V)=F$, with equality of domains and values
at each component.
\end{proof}

\subsection{The recursion category and its splitting}

\begin{proposition}\label{prop:fork-recursion}
Let $\mathscr R_X$ be the full subcategory of $\mathcal A$
whose objects are obtained from $V$ by finitely many applications
of the chosen binary products, binary coproducts, and nonempty-word formation.
Then $\mathscr R_X$ is a recursion category. Its inclusion
into $\mathcal A$ is a range functor, and its restriction
splitting is equivalent to $\mathcal A$.
\end{proposition}
\begin{proof}
The isomorphisms of Lemma~\ref{lem:fork-frames}, applied inductively,
show that every object of $\mathscr R_X$ is isomorphic to $V$.
The binary restriction products, coproducts, and prodominical
zero law restrict from $\mathcal A$ to this full subcategory.
The Turing morphism of Theorem~\ref{thm:fork-turing} has
source $V\times V$ and target $V$ in the subcategory.
Its total parameter maps for objects of $\mathscr R_X$ belong
to $\mathscr R_X$ by fullness. Hence it retains its Turing
universal property there, proving that $\mathscr R_X$ is a
recursion category.

Restrictions and ranges of arrows between objects of $\mathscr R_X$
are endomorphisms of objects of $\mathscr R_X$, so belong to the subcategory.
Their universal properties are unchanged, since all competing
endomorphisms and all arrows between the relevant objects remain
available. Coproducts and their copairings are likewise retained.
Thus the inclusion is a range functor.

For a restriction idempotent $e$ on an object $Y$ of $\mathscr R_X$,
write $m_e:D_e\to Y$ and $r_e:Y\rightharpoonup D_e$ for its
domain inclusion and partial inverse. Send $(Y,e)$ to $D_e$
and an arrow $f:(Y,e)\rightharpoonup(Z,d)$, represented by
$f:Y\rightharpoonup Z$ satisfying $f=dfe$, to
$r_d f m_e:D_e\rightharpoonup D_d$.
The equations $m_er_e=e$ and $r_em_e=1_{D_e}$ show that
identities and composition are preserved. Conversely, send a
partial map $g:D_e\rightharpoonup D_d$ to $m_d g r_e$.
This arrow belongs to $\mathscr R_X$ by fullness and satisfies
$d(m_d g r_e)e=m_d g r_e$. The two assignments are inverse on
each hom-set, so the functor is fully faithful.
It preserves restrictions, since
$\overline{r_d f m_e}=r_e\bar f m_e$.
For every $Z\in\mathcal A$, Lemma~\ref{lem:fork-dominance}
provides a total restriction embedding $m:Z\to V$.
The domain subfunctor of $mm^\dagger$ is isomorphic to $Z$,
so the functor is essentially surjective.
\end{proof}

The same construction gives a functor from the restriction splitting
of $\mathrm{Rec}$ to
$\mathcal K=\operatorname{Par}(\mathcal C_{\mathrm{rec}},
\operatorname{Mono})$: it sends $(X,e)$ to the c.e.\ numerical subset
on which $e$ is the identity, and sends an arrow
$f:(X,e)\rightharpoonup(Y,d)$, represented by a partial recursive
function $f:X\rightharpoonup Y$ satisfying $f=dfe$, to its
restriction between the corresponding subsets. Every c.e.\
numerical carrier embeds by tuple coding into $\N$ and is
isomorphic to the split object of its c.e.\ partial identity
on $\omega$. Partial recursive maps between these carriers are
exactly the partial recursive functions satisfying these equations.
This functor preserves restriction domains and is fully faithful
and essentially surjective. Thus it is a restriction equivalence.

\emph{Decidable equality} on $V$ means a morphism
$\epsilon:V\times V\to1+1$ in $\mathcal E$ whose inverse image of
the first summand is the diagonal subfunctor. Here $1$ is the terminal object of $\mathcal E$.

\begin{proposition}\label{prop:fork-distinction}
The recursion category $\mathscr R_X$ is not equivalent, as a
restriction category, to the classical numerical recursion category.
The object $V$ has no decidable equality in $\mathcal E$.
\end{proposition}
\begin{proof}
The terminal object of $\mathcal E$ is the constant functor at a singleton.
A subobject with empty component $0$ has independently empty or
singleton components $1$ and $2$, giving four subobjects. A
subobject with singleton component $0$ must also have singleton
components $1$ and $2$, giving one more. Thus its subobject order has five elements, whereas
the terminal singleton of $\mathcal C_{\mathrm{rec}}$ has two.
An equivalence of the asserted restriction categories would induce
an equivalence of their restriction splittings, hence of their
total categories $\mathcal E$ and $\mathcal C_{\mathrm{rec}}$.
Equivalences preserve terminal objects and their subobject orders,
so this is impossible.

Suppose that $\epsilon$ decides equality on $V$. The functor $1+1$
is constant and sends both nonidentity morphisms of $J$ to
identity functions. Choose distinct $u,v\in V_0$ differing only
in their $z$ coordinate. Then $r_1^Vu=r_1^Vv$, so
$\epsilon_1(r_1^Vu,r_1^Vv)$ lies in the first summand.
Naturality forces $\epsilon_0(u,v)$ to lie in the same summand,
contrary to $u\ne v$ and the prescribed diagonal inverse image.
\end{proof}

\section{Finite syntax and proofs}
\label{app:fork-proofs}

Let $\mathcal E=\mathcal C_{\mathrm{rec}}^{J}$, where
$J=(1\leftarrow0\to2)$, and let
$\mathcal A=\operatorname{Par}(\mathcal E,\operatorname{Mono})$.
The components of the atom object $V$ are $V_0=\N^4$ and
$V_1=V_2=\N^2$; write $r_i=r_i^V=V(0\to i):V_0\to V_i$
for $i=1,2$.
We construct syntax and finite proofs in $\mathcal E$; their
conclusion image will determine a restriction idempotent in $\mathcal A$.

\subsection{The syntax algebra}

Fix recursive encodings of finite rooted ordered trees and of finite
tuples of natural numbers, with recursive parsing operations. For
$j\in\{0,1,2\}$, let $S_j$ be the set of finite formula trees
generated by
\[
 \theta_j(x)\quad(x\in V_j),\qquad
 \nu_j(\varphi)\quad(\varphi\in S_j),\qquad
 \varphi\Rightarrow_j\psi\quad(\varphi,\psi\in S_j).
\]
The three constructor tags are distinct. Tree parsing and the
decidable numerical carriers $V_j$ make $S_j$ a recursive set of
codes. For $i=1,2$, define $\sigma_i:S_0\to S_i$ by replacing
each atom label $x$ with $r_i(x)$ and retaining every constructor.
These total recursive functions define the functor
$S:J\to\mathcal C_{\mathrm{rec}}$ by $S(j)=S_j$ and
$S(0\to i)=\sigma_i$. The constructors form natural transformations
$\theta:V\to S$, $\nu:S\to S$, and
$\Rightarrow:S\times S\to S$, hence morphisms of $\mathcal E$.

The covariant functor $P:\mathcal E\to\mathcal E$ is defined by
\[
 P(Z)=V+Z+(Z\times Z),\qquad
 P(f)=1_V+f+(f\times f)
 \quad(f:Z\to Z'\text{ in }\mathcal E).
\]
Its identity and composition laws follow componentwise from those
of tagged sums and ordered pairs. The constructor morphisms define
$\iota=[\theta,\nu,\Rightarrow]:P(S)\to S$.

\begin{proposition}\label{prop:fork-syntax-initial}
The algebra $\iota:P(S)\to S$ is initial in the category of
$P$-algebras in $\mathcal E$: for every morphism
$\alpha:P(Z)\to Z$ of $\mathcal E$, there is a unique morphism
$u:S\to Z$ of $\mathcal E$ satisfying $u\iota=\alpha P(u)$.
\end{proposition}
\begin{proof}
At component $j$, define $u_j:S_j\to Z_j$ by evaluating a tree
from its leaves to its root, using the composites of
$\alpha_j:V_j+Z_j+Z_j^2\to Z_j$ with the three coproduct injections.
Induction on tree size proves termination and membership
of the output in $Z_j$. The evaluation is partial recursive on
numerical codes and total on $S_j$, as required for a morphism of
$\mathcal C_{\mathrm{rec}}$.

Write $z_i=Z(0\to i):Z_0\to Z_i$ for $i=1,2$. Naturality of
$\alpha$ says that applying $z_i$ after one constructor operation
is the same as first applying $r_i$ to an atom label, or $z_i$
to each recursive argument, and then using the corresponding
operation at component $i$. Induction on the tree consequently
gives $z_i u_0=u_i\sigma_i$. Thus the component maps form
$u:S\to Z$ in $\mathcal E$, and their recursive definition gives
the algebra equation. Any morphism satisfying that equation agrees
with $u$ on atoms and then on negations and implications by the
same finite induction, proving uniqueness.
\end{proof}

\subsection{The proof object and its range}

Let $a,b$ be disjoint restriction idempotents on $V$ in $\mathcal A$.
Let $A,B\subseteq V$ be the c.e.\ subfunctors on which
$a$ and $b$ are the partial identities, with component subsets
$A_j,B_j\subseteq V_j$. Disjointness
means $A_j\cap B_j=\varnothing$ at each component.
For formulas $\alpha,\beta,\gamma\in S_j$, put
\begin{align*}
 H_{1,j}(\alpha,\beta)
   &=\alpha\Rightarrow_j(\beta\Rightarrow_j\alpha),\\
 H_{2,j}(\alpha,\beta,\gamma)
   &=(\alpha\Rightarrow_j(\beta\Rightarrow_j\gamma))
     \Rightarrow_j((\alpha\Rightarrow_j\beta)
                          \Rightarrow_j(\alpha\Rightarrow_j\gamma)),\\
 H_{3,j}(\alpha,\beta)
   &=(\nu_j\beta\Rightarrow_j\nu_j\alpha)
                          \Rightarrow_j(\alpha\Rightarrow_j\beta).
\end{align*}
A proof at component $j$ is a finite rooted ordered tree with
the following node data and conclusion. A leaf is either a
tagged logical instance $H_{1,j}(\alpha,\beta)$,
$H_{2,j}(\alpha,\beta,\gamma)$, or
$H_{3,j}(\alpha,\beta)$, carrying its displayed formula
parameters; a positive leaf labeled by $x\in A_j$, with
conclusion $\theta_j(x)$; or a negative leaf labeled by
$x\in B_j$, with conclusion $\nu_j\theta_j(x)$.
A nonleaf has the modus-ponens tag and two ordered children.
If their conclusions are $\alpha$ and
$\alpha\Rightarrow_j\beta$, its conclusion is $\beta$.

Let $\mathrm{Prf}_j$ be the set of codes of these trees.
Tree shape, formula syntax, and premise matching can be checked
by recursive operations on the numerical codes. Membership tests
for positive and negative leaves are c.e.; there are only finitely
many such tests in any tree. Dovetailing these finite collections
of tests over all tree codes enumerates exactly $\mathrm{Prf}_j$.
Hence $\mathrm{Prf}_j$ is a c.e.\ carrier.
Reading conclusions from leaves to root gives an arrow
$c_j:\mathrm{Prf}_j\to S_j$ of $\mathcal C_{\mathrm{rec}}$.

For $i=1,2$, replace formula parameters by their images under
$\sigma_i$, positive or negative labels $x$ by $r_i(x)$, and
retain every node tag. The resulting map
$\rho_i:\mathrm{Prf}_0\to\mathrm{Prf}_i$ is an arrow of
$\mathcal C_{\mathrm{rec}}$.
It preserves validity: the inclusions $r_i(A_0)\subseteq A_i$
and $r_i(B_0)\subseteq B_i$ handle signed leaves, and
$\sigma_i$ commutes with the formula constructors and hence
preserves the equalities required at modus-ponens nodes.
Induction on proofs gives
\[
 c_i\rho_i=\sigma_i c_0.
\]
The proof carriers and the morphisms $\rho_i$ define the functor
$\mathrm{Prf}:J\to\mathcal C_{\mathrm{rec}}$ by
$\mathrm{Prf}(j)=\mathrm{Prf}_j$ and
$\mathrm{Prf}(0\to i)=\rho_i$. Their conclusion morphisms form a
natural transformation $c:\mathrm{Prf}\to S$, hence a morphism of
$\mathcal E$.

\begin{theorem}\label{thm:fork-proof-presentation}
For disjoint restriction idempotents $a,b$ on $V$ in $\mathcal A$, the syntax and
proof objects above define a restriction idempotent $t_0=\ran(c)$ on $S$ in
$\mathcal A$. The subfunctor $D_{t_0}\subseteq S$ is the least c.e.\
subfunctor containing all instances of the three logical schemes and the
positive and negative atomic axioms, and closed under modus ponens.
Moreover,
\[
 a\le\theta^*t_0,\qquad
 b\le\theta^*\nu^*t_0,\qquad
 t_0\,\nu^*t_0=0_S.
\]
At each component, these conclusions form a consistent theory of
classical propositional logic.
\end{theorem}
\begin{proof}
Let $T_j=c_j(\mathrm{Prf}_j)\subseteq S_j$.
These sets are c.e.\ images of c.e.\ carriers under
arrows of $\mathcal C_{\mathrm{rec}}$. The conclusion equation gives
$\sigma_i(T_0)\subseteq T_i$, so they form a c.e.\
subfunctor $T\subseteq S$. Its partial identity is
$\ran(c)$: it receives $c$, and the subfunctor of every restriction
idempotent receiving $c$ contains its component images. Thus $T=D_{t_0}$.

Each leaf supplies a proof of its axiom instance. Given proofs
of $\alpha$ and $\alpha\Rightarrow_j\beta$, adjoining a
modus-ponens root supplies a proof of $\beta$. Thus $T$ is
closed under modus ponens. If another c.e.\ subfunctor contains all
the axiom instances and is closed under modus ponens, induction
on each proof puts its conclusion in that subfunctor. This
proves leastness. The one-node positive and negative proofs define
morphisms $A\to\mathrm{Prf}$ and $B\to\mathrm{Prf}$ of $\mathcal E$.
Their composites with $c$ are $\theta|_A$ and
$\nu\theta|_B$, respectively, proving the first two inequalities.

For consistency, fix a component $j$ and suppose two proofs
$p,q\in\mathrm{Prf}_j$ conclude $\varphi$ and $\nu_j\varphi$.
Let $L\subseteq A_j$ be the finite set of atom labels at
positive leaves of either proof. Assign truth value $1$ to an
atom precisely when its label belongs to $L$, and assign $0$
to every other atom. Extend these values to formulas using
Boolean negation and implication.
All positive leaves of both proofs are true. Each negative
leaf has a label in $B_j$, which is disjoint from $A_j$ and
hence from $L$, so that leaf is true as well.

The three logical schemes are Boolean tautologies, and modus
ponens preserves truth. Induction on
$p$ and $q$ would therefore make $\varphi$ and its negation
both true, a contradiction.

Consequently $T_j\cap\nu_j^{-1}(T_j)=\varnothing$ for every
$j$. Restriction idempotents, inverse images, and their meets are componentwise
in $\mathcal A$, so this is precisely $t_0\nu^*t_0=0_S$.
\end{proof}

\subsection{The universal property of the proof object}\label{app:proof-initial}

The conclusion morphism $c:\mathrm{Prf}\to S$ is an object of the
slice category $\mathcal E/S$. We characterize it by the operations
that form axiom leaves and apply modus ponens.

The logical schemes define morphisms
$H_1,H_3:S\times S\to S$ and $H_2:S^3\to S$ of $\mathcal E$.
Let
\[
 \mathrm{Ax}=(S\times S)+S^3+(S\times S)+A+B,
 \qquad
 \lambda=[H_1,H_2,H_3,\theta|_A,\nu\theta|_B]:
       \mathrm{Ax}\longrightarrow S.
\]
An element of $\mathrm{Ax}_j$ records the tag and parameters of an
axiom instance; $\lambda_j$ gives its conclusion.
Write $\pi_1,\pi_2:S\times S\to S$ for the projections and put
\[
 \mu=\langle\pi_1,\Rightarrow\rangle:S\times S\longrightarrow S\times S.
\]
Thus $\mu_j(\alpha,\beta)=(\alpha,\alpha\Rightarrow_j\beta)$.

For an object $q:Q\to S$ of $\mathcal E/S$, form the pullback
\[
\begin{tikzcd}[column sep=large,row sep=large]
 M(Q,q) \arrow[r,"v_q"] \arrow[d,"w_q"'] &
 Q\times Q \arrow[d,"q\times q"] \\
 S\times S \arrow[r,"\mu"'] & S\times S.
\end{tikzcd}
\]
At component $j$, its elements are tuples
$(p_1,p_2,\alpha,\beta)$ satisfying
$q_j(p_1)=\alpha$ and
$q_j(p_2)=\alpha\Rightarrow_j\beta$.
The morphism
$\kappa_q=\pi_2w_q:M(Q,q)\to S$ gives the conclusion $\beta$.

Define an endofunctor $\mathsf H:\mathcal E/S\to\mathcal E/S$ by
\[
 \mathsf H(Q,q)=
 \bigl(\mathrm{Ax}+M(Q,q),[\lambda,\kappa_q]\bigr).
\]
For a slice morphism $h:(Q,q)\to(R,r)$, so that $rh=q$,
the pullback property gives a unique
$M(h):M(Q,q)\to M(R,r)$ satisfying
\[
 v_rM(h)=(h\times h)v_q,\qquad w_rM(h)=w_q.
\]
Put $\mathsf H(h)=1_{\mathrm{Ax}}+M(h)$.
The second equation gives $\kappa_rM(h)=\kappa_q$, so this is a
slice morphism. Uniqueness in the pullback proves the identity and
composition laws for $\mathsf H$.

Axiom-leaf formation gives $\ell:\mathrm{Ax}\to\mathrm{Prf}$.
Joining two proofs under a modus-ponens root gives
$m:M(\mathrm{Prf},c)\to\mathrm{Prf}$.
The implication constructor is injective at each component, so the two
children of a valid modus-ponens node determine its consequent uniquely.
The maps $\ell$ and $m$ are morphisms of $\mathcal E$: the operations on tree codes
are recursive, and relabeling commutes with their constructors.
They satisfy $c\ell=\lambda$ and $cm=\kappa_c$.
Consequently
\[
 \delta=[\ell,m]:\mathsf H(\mathrm{Prf},c)
                     \longrightarrow(\mathrm{Prf},c)
\]
is an $\mathsf H$-algebra in $\mathcal E/S$.

\begin{proposition}\label{prop:fork-proof-initial}
The algebra $\delta$ is initial. For every
$\mathsf H$-algebra
$\eta:\mathsf H(Q,q)\to(Q,q)$ in $\mathcal E/S$, there is a unique
morphism $u:\mathrm{Prf}\to Q$ in $\mathcal E$ such that
\[
 qu=c,\qquad
 u\delta=\eta\bigl(1_{\mathrm{Ax}}+M(u)\bigr).
\]
Equivalently, the following square commutes in $\mathcal E/S$:
\[
\begin{tikzcd}[column sep=large,row sep=large]
 \mathsf H(\mathrm{Prf},c) \arrow[r,"\delta"]
       \arrow[d,"\mathsf H(u)"'] &
 (\mathrm{Prf},c) \arrow[d,"u"] \\
 \mathsf H(Q,q) \arrow[r,"\eta"'] & (Q,q).
\end{tikzcd}
\]
\end{proposition}
\begin{proof}
At component $j$, evaluate a finite proof tree from its leaves to
its root. At an axiom leaf with datum $\xi\in\mathrm{Ax}_j$, apply
$\eta_j$ to the corresponding element of the first summand.
Since $\eta$ is a slice morphism, the result has conclusion
$\lambda_j(\xi)$.

At a modus-ponens node, suppose the two subtrees conclude
$\alpha$ and $\alpha\Rightarrow_j\beta$.
By induction, their values $p_1,p_2\in Q_j$ have those same
conclusions under $q_j$. Hence
$(p_1,p_2,\alpha,\beta)$ belongs to $M(Q,q)_j$.
Apply $\eta_j$ to its image in the second summand. The result
has conclusion $\beta$.

This recursion defines $u_j:\mathrm{Prf}_j\to Q_j$ and proves
$q_ju_j=c_j$. Each step uses a partial recursive function total
on its specified carrier. The tree is finite, so its evaluation
is total on $\mathrm{Prf}_j$ and is an arrow of
$\mathcal C_{\mathrm{rec}}$.
Naturality of $\eta$, $\lambda$, and the pullback projections
shows inductively that evaluation commutes with relabeling.
Thus the maps $u_j$ form a morphism $u:\mathrm{Prf}\to Q$ of
$\mathcal E$. Their defining recursion is the algebra equation.
Any other algebra morphism agrees with $u$ on leaves and then
at each modus-ponens node by induction on the tree, proving uniqueness.
\end{proof}

The fibers of $c_j$ consist of derivation trees with a common
conclusion. Distinct trees remain distinct elements of
$\mathrm{Prf}_j$. The range $\ran(c)$ records only which formulas
have at least one derivation.


\subsection{First-order derivations and sentence proofs}
\label{app:first-order-proofs}

The numerical proof objects in Sections~\ref{sec:miller} and
\ref{sec:arithmetic} have an analogous characterization.
Fix a first-order theory $T$. Let
$\mathcal D$ be $\mathcal C_{\mathrm{rec}}$ for a c.e.\ axiom set,
or $\operatorname{Tot}(\operatorname{Split}(\mathcal E_a))$ for an
$\mathcal E_a$-semidecidable axiom set. Use a sound and complete
Hilbert calculus with a recursive set of logical axioms and the rules
of modus ponens and generalization. The nonlogical axioms are sentences.
Write $\mathsf F$ for the recursive carrier of all formulas,
$\mathsf S\hookrightarrow\mathsf F$ for its sentence subobject, and
$\mathsf L,\mathsf T\hookrightarrow\mathsf F$ for the logical and
nonlogical axiom carriers.

The objects of rule instances are
\[
 R_0=\mathsf L+\mathsf T,\qquad
 R_1=\N\times\mathsf F,\qquad R_2=\mathsf F\times\mathsf F.
\]
For $k=0,1,2$, let $p_k:R_k\to\mathsf F^k$ and
$q_k:R_k\to\mathsf F$ give the premises and conclusion. Here
$\mathsf F^0=1$, $p_0$ is the terminal map, and $q_0$ is the copairing
of the two axiom inclusions. For generalization and modus ponens, put
\[
\begin{aligned}
 p_1(v,\varphi)&=\varphi,&
 q_1(v,\varphi)&=\forall v\,\varphi,\\
 p_2(\varphi,\psi)&=(\varphi,\varphi\Rightarrow\psi),&
 q_2(\varphi,\psi)&=\psi.
\end{aligned}
\]
Variables are indexed by $v\in\N$. Generalization has no additional
side condition because all nonlogical assumptions are sentences.
The displayed connectives are constructors on formula codes.
These maps are morphisms of $\mathcal D$.

For $(Q,q)\in\mathcal D/\mathsf F$, put
\[
 B_k(Q,q)=R_k\times_{\mathsf F^k}Q^k,
 \qquad
 \mathsf H_T(Q,q)=
 \left(\coprod_{k=0}^2 B_k(Q,q),
       [q_k\operatorname{pr}_{R_k}]_{k=0}^2\right).
\]
The pullback uses $p_k$ and $q^k$; its elements are a rule instance
and proofs of its premises. A slice morphism $h$ induces
$B_k(h)$ by applying $h$ to each premise, and
$\mathsf H_T(h)=\coprod_k B_k(h)$.
Pullback uniqueness gives the identity and composition laws.
The displayed products, pullbacks, and coproducts exist in
$\mathcal D$: their carriers are obtained by finite products,
equality tests, and tagged unions of the specified semidecidable sets.

\begin{proposition}\label{prop:first-order-proof-initial}
Let $\mathsf{Der}_T$ consist of finite trees labeled by these rule
instances, with matching premise conclusions, and let
$d_T:\mathsf{Der}_T\to\mathsf F$ give their root conclusions.
Adjoining a rule at the root defines an algebra
\[
 \delta_T:\mathsf H_T(\mathsf{Der}_T,d_T)
       \longrightarrow(\mathsf{Der}_T,d_T).
\]
This algebra is initial in $\mathcal D/\mathsf F$.
The proof object with sentence conclusions is the pullback
\[
\begin{tikzcd}[column sep=large,row sep=large]
 \mathrm{Prf}_T \arrow[r] \arrow[d,"c_T"'] &
 \mathsf{Der}_T \arrow[d,"d_T"] \\
 \mathsf S \arrow[r,hook] & \mathsf F.
\end{tikzcd}
\]
Viewed in $\operatorname{Par}(\mathcal C_{\mathrm{rec}},\operatorname{Mono})$
or $\operatorname{Split}(\mathcal E_a)$, respectively, the map $c_T$
has range equal to the theorem restriction idempotent on $\mathsf S$.
\end{proposition}
\begin{proof}
Finite tree shape, rule parameters, and premise matching are recursive
tests. The finitely many nonlogical axiom tests are semidecidable in
the stated sense, so the valid trees form an object of $\mathcal D$.
The conclusion map and the operation adjoining a rule at the root are
recursive on valid trees and satisfy the conclusion equations.

For any $\mathsf H_T$-algebra, evaluate a tree from its leaves to
its root using its algebra operations. Induction on the tree proves
termination, preservation of the conclusion, and the algebra-morphism
equation. The same induction proves uniqueness. In
$\operatorname{Tot}(\operatorname{Split}(\mathcal E_a))$, first run the
membership test for the valid-tree carrier. Evaluation then has exactly
that domain; its value is independent of the name because each algebra
operation has this property. Thus the unique evaluation map belongs to
$\mathcal D$ in both cases.

The pullback selects precisely the trees whose root conclusions are
sentences. Closedness is recursive on formula codes, so this selection
preserves semidecidability. Its image under the conclusion map is the
set of theorems.
\end{proof}

\section{Necessity of section lifting for inseparability}
\label{app:lifting-counterexample}

The lifting hypothesis in Theorem~\ref{thm:local-lifting} cannot be
omitted, even for a faithful range functor that preserves restriction
products and reflects zero arrows.

Fix an acceptable numbering $(\varphi_e)_{e\in\N}$ of the ordinary
partial recursive functions $\N\rightharpoonup\N$. For a partial
numerical function $f$, write $f(x)\mathord\downarrow$ when $f$ is
defined at $x$. Put
\[
 K=\{e:\varphi_e(e)\mathord\downarrow\}.
\]
The set $K$ is not recursive. Indeed, a decision procedure for $K$
would give a program that, on input $x$, halts precisely when the
procedure answers that $x\notin K$. Applying this program to its own
index gives a contradiction.
Let $(W_a^K)_{a\in\N}$ enumerate the domains of an acceptable
numbering of partial functions recursive relative to $K$.
Fix a recursive bijective pairing function
$\langle-,-\rangle:\N^2\to\N$, with projections $\pi_0,\pi_1$.

The objects of $\mathcal C_K$ are the formal numerical types generated
from $\N$ by binary products and tagged coproducts, with fixed recursive
codings. An arrow $f:X\rightharpoonup Y$ is a partial function whose
domain $D_f$ is c.e.\ relative to $K$ and which has an ordinary partial
recursive extension $f_0:X\rightharpoonup Y$. Thus $f_0$ is defined
on every element of $D_f$ and agrees with $f$ there. Equality is equality
of partial-function graphs. These arrows need not themselves be ordinary
partial recursive functions. The total arrows of $\mathcal C_K$ are exactly
the ordinary total recursive functions between the indicated types.

On $\N$, the arrows are precisely the functions
\[
 f_{e,a}(x)\simeq
 \begin{cases}
  \varphi_e(x),&x\in W_a^K,\\
  \text{undefined},&x\notin W_a^K.
 \end{cases}
\]
Here $\simeq$ includes simultaneous undefinedness: the actual domain
is $W_a^K\cap\operatorname{dom}(\varphi_e)$. Every arrow has this
form by taking an index of its extension and an oracle index of its
specified domain.

\begin{proposition}\label{prop:missing-lift}
The category $\mathcal C_K$ is a prodominical recursion category with
binary restriction coproducts and ranges. Its inclusion
$F:\mathcal C_K\to\Pfn$ satisfies every hypothesis of
Theorem~\ref{thm:local-lifting} except the lifting hypothesis.
For the Turing morphism and coproduct isomorphism constructed below,
the specified pair is separated by complementary restriction idempotents in
$\mathcal C_K$. Moreover, there are parallel arrows in $\mathcal C_K$
with disjoint restriction idempotents and no common upper bound in
the restriction order, although finite joins of restriction idempotents
exist.
\end{proposition}

\begin{proof}
Let $f:X\rightharpoonup Y$ and $g:Y\rightharpoonup Z$ have
ordinary partial recursive extensions $f_0,g_0$. Their composite has
extension $g_0f_0$ and domain
\[
 D_{gf}=D_f\cap f_0^{-1}(D_g).
\]
This domain is c.e.\ relative to $K$: enumerate $D_f$, compute the
corresponding $f_0$-values, and test for their appearance in the
enumeration of $D_g$, using dovetailing. Identities and associativity
are inherited from partial functions. Restriction is the partial
identity on $D_f$, which has the ordinary identity as an extension.
Empty partial functions give restriction zeros.

Cartesian products and tagged copairings preserve the arrow class:
their domains are respectively products and tagged unions of the given
domains, and their extensions are the corresponding ordinary partial
recursive programs. Consequently the set-theoretic restriction
products and coproducts have their usual universal properties in
$\mathcal C_K$. Product with a zero arrow is zero, and the chosen
recursive coding isomorphisms make the category isotypical.

The set $f(D_f)$ is c.e.\ relative to $K$: enumerate $D_f$ and run
$f_0$ on its elements. The partial identity on this image belongs to
$\mathcal C_K$ and is its range. Leastness and the agreement law
hold for partial functions: equality after $f$ gives equality on
every value in its image, including equality of definedness.
Ranges commute with restriction products because images of products
of partial functions are products of images. Restriction idempotents
are the partial identities on the subsets c.e.\ relative to $K$;
their order corresponds to inclusion, and their finite joins are
partial identities on unions.

Define $\tau_K:\N\times\N\rightharpoonup\N$ by
\[
 \tau_K(\langle e,a\rangle,x)\simeq f_{e,a}(x).
\]
Its domain is c.e.\ relative to $K$, uniformly in all three
variables, and it has the ordinary partial recursive extension
$(\langle e,a\rangle,x)\mapsto\varphi_e(x)$.
To prove its Turing property, let
$v:\N\times\N\rightharpoonup\N$ be an arrow, with ordinary
partial recursive extension $v_0:\N\times\N\rightharpoonup\N$ and
domain $D_v$ c.e.\ relative to $K$. Specializing the ordinary program
for $v_0$ gives an ordinary total recursive function $e:\N\to\N$ with
\[
 \varphi_{e(z)}(x)\simeq v_0(z,x).
\]
Specializing an oracle program for $D_v$ gives an ordinary total
recursive function $a:\N\to\N$ with
\[
 W_{a(z)}^K=\{x:(z,x)\in D_v\}.
\]
The latter specialization is ordinary recursive because it only
inserts the finite numeral for $z$ into program code; the oracle is
used when the resulting program runs. Thus the function
$r:\N\to\N$ given by $r(z)=\langle e(z),a(z)\rangle$ is a total
morphism of $\mathcal C_K$, and
$v=\tau_K(r\times1_{\N})$, with equality of domains and values.
The fixed recursive codings give the same argument for every formal
parameter type. Hence $\tau_K$ is a Turing morphism.

The inclusion $F$ is faithful and preserves restrictions, zeros,
ranges, binary restriction products, and coproducts. It also reflects
zero arrows. Its target $\Pfn$ is split and locally connected:
restriction idempotents split through their subsets of definition,
and partial identities on singletons detect partial functions.
Partial injections have partial inverses, giving weak choice, and
inverse images preserve unions. Also $F(1_{\N})\ne0$.

Use the total recursive isomorphism
\[
 b:\N\longrightarrow\N+\N,\qquad
 b(2n)=i_0(n),\quad b(2n+1)=i_1(n).
\]
Put $d=\tau_K\Delta_{\N}$ and
$a_0=(bd)^*(1_{\N}+0_{\N})$,
$a_1=(bd)^*(0_{\N}+1_{\N})$.
For $q\in\N$, write
$p(q)\simeq\varphi_{\pi_0(q)}(q)$. The carriers of these restriction idempotents are
\[
 \begin{aligned}
 D_{a_0}&=\{q:q\in W_{\pi_1(q)}^K,
                    \ p(q)\mathord\downarrow\text{ with an even value}\},\\
 D_{a_1}&=\{q:q\in W_{\pi_1(q)}^K,
                    \ p(q)\mathord\downarrow\text{ with an odd value}\}.
 \end{aligned}
\]
Let
\[
 A=\{q:p(q)\mathord\downarrow\text{ with an even value}\}.
\]
This is an ordinary c.e.\ set and is decidable relative to $K$.
For the latter assertion, uniformly construct a program that ignores
its input, simulates $p(q)$, and halts precisely when the output is
even; membership of its index in $K$ decides membership of $q$ in $A$.
Thus $A$ and $A^c$ are c.e.\ relative to $K$, and their partial
identities are complementary restriction idempotents in $\mathcal C_K$.
They separate the displayed pair because
$D_{a_0}\subseteq A$ and $D_{a_1}\subseteq A^c$.

The set $A$ is not recursive. Fix any index $a_*$. Choose an ordinary
total recursive function $h:\N\to\N$ such that $\varphi_{h(n)}$
ignores its input, simulates $\varphi_n(n)$, and returns $0$ if that
simulation halts. Then
\[
 n\in K\quad\Longleftrightarrow\quad
 \langle h(n),a_*\rangle\in A.
\]
Write $1_B$ for the partial identity on a subset $B$. The source
arrow $p_A=\nabla_{\N}(1_{A^c}+1_A)$ is sent to the partial
injection whose unique section in $\Pfn$ is
\[
 s_A(n)=
 \begin{cases}
  i_1(n),&n\in A,\\
  i_0(n),&n\notin A.
 \end{cases}
\]
This section is total and not recursive, since its tag decides $A$.
If it had a preimage under the inclusion $F$, that source arrow would
be total and hence ordinary total recursive, a contradiction.
This is the missing lift required by
Theorem~\ref{thm:local-lifting} for the complemented restriction idempotent $1_A$.

Finally, the arrows $i_1 1_A$ and $i_0 1_{A^c}$ belong to
$\mathcal C_K$: they have the total recursive extensions $i_1$
and $i_0$, respectively. Their domains are disjoint. Any common
upper bound would agree with the first arrow on $A$ and the second
on $A^c$. It would therefore be the total function $s_A$, which
is not a source arrow. Thus no such upper bound exists.
\end{proof}

\begin{thebibliography}{99}
\small
\bibitem{BauerHanson}
A.~Bauer and J.~E. Hanson, \emph{The Countable Reals},
\href{https://arxiv.org/abs/2404.01256v3}{arXiv:2404.01256v3}, 2026.
\bibitem{Bernardi}
C.~Bernardi, \emph{Precomplete equivalence relations in dominical
categories}, Proc. Amer. Math. Soc. \textbf{98} (1986), 124--134.
\bibitem{CockettHofstra}
J.~R.~B. Cockett and P.~J.~W. Hofstra,
\emph{Introduction to Turing categories},
Ann. Pure Appl. Logic \textbf{156} (2008), 183--209.
\bibitem{CockettLack}
J.~R.~B. Cockett and S.~Lack, \emph{Restriction categories. I.
Categories of partial maps}, Theoret. Comput. Sci. \textbf{270}
(2002), 223--259.
\bibitem{CockettLackIII}
J.~R.~B. Cockett and S.~Lack,
\emph{Restriction categories III: colimits, partial limits and extensivity},
Math. Structures Comput. Sci. \textbf{17} (2007), 775--817.
\bibitem{CurryCosheaves}
J.~M. Curry, \emph{Dualities between cellular sheaves and cosheaves},
J. Pure Appl. Algebra \textbf{222} (2018), 966--993.
\href{https://doi.org/10.1016/j.jpaa.2017.06.001}
{doi:10.1016/j.jpaa.2017.06.001}.
\bibitem{DPH}
R.~A. Di Paola and A.~Heller, \emph{Dominical categories:
recursion theory without elements}, J. Symbolic Logic
\textbf{52} (1987), 594--635.
\bibitem{DPM}
R.~A. Di Paola and F.~Montagna, \emph{Some properties of the syntactic
$p$-recursion categories generated by consistent, recursively enumerable
extensions of Peano arithmetic}, J. Symbolic Logic \textbf{56} (1991),
643--660.
\bibitem{Heller}
A.~Heller, \emph{An existence theorem for recursion categories},
J. Symbolic Logic \textbf{55} (1990), 1252--1268.
\bibitem{Lengyel2004}
F.~Lengyel, \emph{More existence theorems for recursion categories},
Ann. Pure Appl. Logic \textbf{125} (2004), 1--41.
\bibitem{Corrections}
F.~Lengyel, \emph{Corrections to More existence theorems for recursion
categories}, version 10, 20 October 2026,
\href{https://github.com/flengyel/errata/blob/5555e4fa1da7b8e3a0836dd1fbf1ca385cacd03b/lengyel2004/lengyel2004_repairs_v10.tex}{source}.
\bibitem{Composition}
F.~Lengyel, \emph{Pointwise provable equality and the failure of
composition}, 2026,
\href{https://arxiv.org/abs/2609.25556}{arXiv:2609.25556}.
\bibitem{Miller2004}
J.~S. Miller, \emph{Degrees of unsolvability of continuous functions},
J. Symbolic Logic \textbf{69} (2004), no.~2, 555--584.
\href{https://doi.org/10.2178/jsl/1082418543}{doi:10.2178/jsl/1082418543}.
\bibitem{Montagna}
F.~Montagna, \emph{``Pathologies'' in two syntactic categories of partial
maps}, Notre Dame J. Formal Logic \textbf{30} (1989), 105--116.
\bibitem{RosoliniThesis}
G.~Rosolini, \emph{Continuity and effectiveness in topoi}, D.Phil.
thesis, University of Oxford, 1986.
\bibitem{Rosolini}
G.~Rosolini, \emph{Representation theorems for $P$-categories},
in F.~Borceux (ed.), \emph{Categorical Algebra and its Applications},
Lecture Notes in Mathematics \textbf{1348}, Springer, Berlin, 1988,
307--315.
\end{thebibliography}
\end{document}

